%%
%% This file follows the user's own elsarticle-template-num.tex, uploaded
%% for this project. Preamble, proof formatting, and frontmatter structure
%% are matched to that template; content (title, abstract, keywords, body)
%% is this paper's own.
\documentclass[final,5p,twocolumn]{elsarticle}

\usepackage{amsmath,amsfonts}
\usepackage{algorithm}
\usepackage{algpseudocode}
\renewcommand{\Comment}[1]{\hfill$\triangleright$ #1}
\usepackage{array}
\usepackage[font=normalsize,labelfont=bf]{caption}
\usepackage[]{subcaption}
\usepackage{arydshln}
\usepackage{textcomp}
\usepackage{stfloats}
\usepackage{url}
\usepackage{verbatim}
\usepackage{graphicx}
\usepackage{IEEEtrantools}
\usepackage{hyperref}
\usepackage{setspace}
\usepackage{soul}
\usepackage{xcolor}
\usepackage{bm}
\usepackage{amsthm}
\hyphenation{op-tical net-works semi-conduc-tor IEEE-Xplore}
\def\BibTeX{{\rm B\kern-.05em{\sc i\kern-.025em b}\kern-.08em
		T\kern-.1667em\lower.7ex\hbox{E}\kern-.125emX}}
\usepackage{balance}
\setlength{\textfloatsep}{8pt}
\newcommand{\orcidlink}[1]{\href{https://orcid.org/#1}{\includegraphics[scale=0.06]{pictures/orcid.png}}}

\newtheorem{definition}{Definition}
\newtheorem{lemma}{Lemma}
% The template defines Definition and Lemma only; Corollary and Remark are
% this paper's own addition, chained to the same counter as Lemma so all
% three number in one sequence through the appendix.
\newtheorem{corollary}[lemma]{Corollary}
\theoremstyle{remark}
\newtheorem{remark}{Remark}
\theoremstyle{plain}

\raggedbottom
\makeatletter
\usepackage{xcolor,etoolbox}
\makeatletter
\renewcommand\proofname{\textcolor{black}{\textbf{Proof}}}

\patchcmd{\proof}
{\item[\hskip\labelsep\scshape #1\@addpunct{.}]}
{\item[\hskip\labelsep\bfseries\color{black}#1\@addpunct{:}] \mbox{}\\}
{}{}%

\patchcmd{\proof}
{\item[\hskip\labelsep #1\@addpunct{.}]}
{\item[\hskip\labelsep\bfseries\color{black}#1\@addpunct{:}] \mbox{}\\}
{}{}%

\patchcmd{\proof}
{\item[\hskip\labelsep\textit{#1}\@addpunct{.}]}
{\item[\hskip\labelsep\bfseries\color{black}#1\@addpunct{:}] \mbox{}\\}
{}{}%
\makeatother
\renewcommand{\qedsymbol}{}

\newcommand{\TBD}{\textbf{TBD}} % placeholder for the one confirmed unrun analysis (density sweep, if relevant here); real results go in directly elsewhere

\journal{Journal of Network and Computer Applications}

\begin{document}

\begin{frontmatter}

\title{Byzantine Robust Federated LSTM with STARK Hop Legitimacy Attestation, Witness Consensus Passive Detection, and Controller Failover Zero Trust Blockchain Against Hidden Forwarding Attacks in Software Defined Vehicular Networks}

% TODO: author block pending confirmation of the final author list for this
% specific paper (distinct author ordering may apply versus PHANTOM and
% NEXUS). Single author set below, matching the pattern of the uploaded
% template with its own co-author block commented out.
\author[1]{Patikiri Arachchige Don Shehan Nilmantha Wijesekara\orcidlink{0000-0002-3045-1596}\corref{cor1}\\}
\author[1]{Siyagunakosgodage Apsari Udithara Fernando\orcidlink{0009-0008-1328-8855}}
\ead{nilmantha@eie.ruh.ac.lk}
\cortext[cor1]{Corresponding author: \texttt{nilmantha@eie.ruh.ac.lk}}

\affiliation[1]{organization={Department of Electrical \& Information Engineering, Faculty of Engineering, University of Ruhuna},
            city={Galle},
            postcode={80000},
            country={Sri Lanka}}

\begin{abstract}
	Hidden forwarding attacks against software defined vehicular networks
	leave a packet's original path untouched while a duplicate is secretly
	routed to an unauthorized destination, enabling eavesdropping without
	any observable disruption to normal delivery. Existing path
	validation and flow conservation defenses require enough packets
	observed over a flow's lifetime to test reliably, a requirement a
	vehicle's brief residence inside a single roadside unit's coverage
	zone routinely fails to meet regardless of how the underlying test is
	tuned; separately, cryptographic verification alone cannot distinguish
	a duplicate that is a genuinely authentic retransmission from a
	legitimate second delivery, since authenticity and intent are not the
	same thing. This paper presents VANGUARD HF, spanning signatures
	corrected to require that the accusing node was the addressed
	recipient of a duplicate rather than a bystander that merely
	overheard it, a post quantum proof of next hop legitimacy that
	certifies routing compliance without revealing a node's routing
	choice, a federated classifier trained across roadside units that
	reads a reception log anomaly score no other observable can provide
	for the passive case, and a witness quorum that corroborates an event
	across independently observing vehicles before any trust penalty
	applies. Evaluated across attack penetration, forwarding intensity,
	vehicular speed, network scale, and decreasing observable evidence,
	VANGUARD HF sustains macro Matthews correlation coefficient \TBD\
	against \TBD\ for the one baseline addressing this attack class at
	all. These results show that covering hidden forwarding in full
	requires abandoning the assumption that cryptographic authenticity
	implies legitimate intent, which every reviewed defense still makes.
\end{abstract}

\begin{keyword}
	Software Defined Vehicular Networks \sep Hidden Forwarding Attack \sep Federated Learning \sep Post Quantum Cryptography \sep Byzantine Fault Tolerance \sep Blockchain
\end{keyword}

\end{frontmatter}

\section{Introduction}\label{introduction}

\subsection{Background}

\IEEEPARstart{S}{oftware} defined vehicular networks separate traffic forwarding at
roadside units from routing decisions made by a logically centralized
controller, and this separation creates an attack surface a static
deployment does not have: a compromised entity can duplicate a packet and
route the copy to a destination the sender never authorized, while the
original continues to its intended recipient exactly as expected. Because
nothing about the legitimate delivery is disturbed, the packet's own
sender or intended recipient has no reason to suspect anything happened
at all. This paper addresses this attack class, hidden forwarding, as one
of two disjoint classes this project considers, the other being selective
time delay, developed independently in the companion paper PHANTOM
\cite{Phantom2026}.

The building blocks this attack exploits, a corrupted controller view and
an unverified multi controller forwarding decision, have each been
demonstrated independently before being combined into the threat this
paper addresses. A topology poisoning attack has been carried out
successfully against four mainstream software defined controllers under
realistic vehicular mobility, showing that an attacker can corrupt the
controller's own view of the network directly, which is precisely the
kind of manipulation a hidden forwarding rule depends on to route a
duplicate toward an unauthorized destination without raising suspicion
\cite{Wang2020Topology}. Blockchain based verification that a forwarding
rule is enforced consistently across controller domains has been proposed
separately, using an aggregate signature scheme to check cross domain
correctness \cite{Li2022BlockREV}, though it verifies that a rule was
enforced as issued rather than whether the issued rule itself was
authorized in the first place, which is the distinction Signature S5's
endorsement check in this paper is built around. Early lightweight
vehicular intrusion detection targeted false alert generation directly
\cite{Sedjelmaci2014}, a concern this paper's own addressing conjunct
extends into the signature construction itself rather than treating it as
a separate detection target, and hardware backed authentication using
direct anonymous attestation addressed message privacy before any
centralized controller existed to threaten it \cite{Sumra2011}.

\subsection{Limitations of Existing Work}

Prior defenses close only part of this problem, and each limitation holds
on its own terms regardless of what any particular solution looks like.
Path validation and flow conservation frameworks such as HSA and FADE
test whether a flow's observed forwarding behaviour is internally
consistent over its lifetime, which requires enough packets to
accumulate a statistically reliable sample; a vehicle's residence time
inside a single roadside unit's coverage zone, typically several seconds
to tens of seconds at highway speed, routinely provides fewer packets
than such a test needs, independent of how the test's own parameters are
tuned. Cryptographic accountability schemes such as SDNsec verify that a
packet was forwarded correctly, which a passively retransmitted,
genuinely authentic duplicate also satisfies, since nothing about the
packet's content or signature was altered in producing it; verifying
authenticity and verifying that a delivery was intended are simply not
the same check, and no scheme reviewed in Section~\ref{sec:related}
attempts to distinguish them. Neither limitation depends on the other:
a mobility bounded observation window defeats a flow conservation test
regardless of how cryptographically strong the network's signing scheme
is, and an authenticity check that cannot see past authenticity fails
regardless of how much traffic a flow generates.

\subsection{Problem Statement}

The practical consequence is that a vehicle's location, driving pattern,
or communication content can be exposed to an unauthorized party while
every observable signal available to the sender and the intended
recipient indicates normal operation. A defense that only inspects
whether a flow's statistics look conserved, or only checks that a packet
was signed by its claimed sender, provides no warning in exactly the
passive case where the duplicate is a real, unmodified copy of genuine
traffic, which is not a corner case this paper's threat model sets aside
but the specific case its architecture is built to cover.

\subsection{Research Questions}

This paper is organised around five research questions that the
experiments of Section~\ref{sec:results} answer directly.
\begin{itemize}
	\item \textbf{RQ1:} Does detection quality hold across attack
	penetration and forwarding intensity, and does it differ between
	active variants, where a duplicate is fabricated, and passive
	variants, where it is a genuine retransmission?
	\item \textbf{RQ2:} Does detection quality vary with vehicular speed,
	and if so, is the mechanism responsible timing related, as in the
	companion paper, or specific to how long a witness quorum takes to
	form?
	\item \textbf{RQ3:} Do detection quality and end to end latency remain
	within their required bounds as network scale grows, and where does
	the binding cost actually sit as the network grows along the roadside
	unit axis rather than the vehicle axis?
	\item \textbf{RQ4:} How does detection degrade as attack observable
	evidence decreases along both the selectivity and copy destination
	divergence dimensions specific to this attack class?
	\item \textbf{RQ5:} Across each of the four hidden forwarding variants
	individually, does VANGUARD HF outperform existing defenses, and does
	a defense built for one hidden forwarding mechanism transfer to the
	other?
\end{itemize}

\subsection{Contributions}

The principal contributions of this paper are:
\begin{itemize}
	\item \textbf{Broadcast correct attack modeling:} the formal
	definition of four hidden forwarding variants across the control and
	data plane, corrected to require that the accusing node was the
	addressed recipient of a duplicate, closing a construction flaw that
	otherwise lets any bystander that merely overheard a broadcast
	transmission satisfy an accusation condition meant for a genuine
	recipient.
	\item \textbf{Post quantum hop legitimacy attestation:} a zero
	knowledge proof that a forwarding node's next hop choice complies with
	its authorized routing policy, certifying routing compliance without
	revealing the routing choice itself to any verifier.
	\item \textbf{Federated classification with a duplication specific
	observable:} a classifier trained across roadside units without
	centralizing raw traffic, reading a reception log anomaly score that
	remains zero under no attack and rises under every hidden forwarding
	variant including the passive ones, where every cryptographic
	verification available otherwise passes.
	\item \textbf{Witness quorum corroboration:} a mechanism requiring
	independent confirmation from multiple observing vehicles before any
	trust penalty applies, covering the passive case no cryptographic or
	statistical observable can resolve alone.
	\item \textbf{Byzantine robust federated aggregation:} an aggregation
	procedure that excludes a roadside unit's model update once its
	submitted weights deviate beyond a bounded distance from the
	consensus, preventing a single compromised roadside unit from
	corrupting the classifier every other roadside unit relies on.
\end{itemize}

\subsection{Novelty}

This is the first defense against hidden forwarding in software defined
vehicular networks to identify and correct a construction flaw present in
every broadcast medium signature scheme reviewed in
Section~\ref{sec:related}: without an explicit addressing check, any
bystander that overheard a transmission satisfies an accusation condition
intended for a genuine recipient, and this paper's correction is a
structural fix to the signature itself, not a refinement of an existing
threshold. It is also the first framework in this space to provide direct
learned signal for passive hidden forwarding independently of any
cryptographic observable, through a reception log anomaly score that no
prior scheme computes, since every reviewed defense either tests
cryptographic authenticity, which a passive duplicate satisfies by
construction, or flow level statistics, which a bounded observation
window undermines. The witness quorum mechanism completes this coverage
for the case that remains after both the corrected signatures and the
classifier are exhausted, corroboration from independent physical
observation rather than from any signal the compromised node itself
produces. Collectively, this paper establishes that passive hidden
forwarding, the case every baseline in Section~\ref{sec:related} misses
entirely, is detectable without relying on any assumption the passive
attack is specifically designed to defeat, a combination this paper's
review of prior work finds nowhere else.

\subsection{Paper Organization}

The remainder of this paper is organised as follows. Section~\ref{sec:related}
reviews related work and positions the gap this paper addresses.
Section~\ref{system_model} formalises the network model, threat model,
and four attack variants. Section~\ref{sec:framework} presents the
VANGUARD HF framework. Sections~\ref{sec:evaluation_setup}
and~\ref{sec:results} describe the evaluation setup and report the five
experiments. Section~\ref{sec:discussion} discusses the results, and the
final section concludes.


\section{Related Work}
\label{sec:related}

Hidden forwarding attacks against software defined vehicular networks leave
the original packet on its legitimate path while a duplicate is silently
routed elsewhere, so the literature relevant to detecting it spans path
validation, forwarding anomaly detection, and cryptographic accountability
for the data plane, three lines of work that each catch part of the problem
without fully closing it.

Path validation was formalised early by Mohammadi et al., who correlate the
forwarding rules a network has installed against the forwarding status
packets actually exhibit, using header space analysis to detect a compromised
switch that drops, copies, or misroutes traffic \cite{Mohammadi2016}. Pang
et al. propose FADE, a flow conservation based detector that flags duplication
or path deviation by testing whether the volume entering and leaving a flow
remains consistent \cite{Pang2016FADE}. Both build on the older idea of
reconstructing a packet's actual path from consistent hash based sampling
across network domains \cite{Duffield2001}, and both share a structural
weakness against a mobile, short lived flow: they need enough packet samples
over a flow's observed lifetime to test conservation reliably, and neither
was designed for a setting where that lifetime is bounded by how long a
vehicle remains inside a single roadside unit's coverage zone.

Cryptographic accountability for the forwarding path itself has been pursued
independently of flow level statistics. Sasaki et al. propose SDNsec, which
embeds a symmetric key cryptographic marking in every packet so that each
switch proves it forwarded the packet correctly, directly addressing the path
deviation attack this paper's active hidden forwarding variants also produce
\cite{Sasaki2016SDNsec}. Bargayary and Medhi propose FTISCON, which anchors flow
rule installation to a smart contract so an installed rule can be checked
against a committed policy after the fact \cite{Bargayary2023FTISCON}. Neither
extends this accountability to a zero knowledge setting: SDNsec's markings are
verifiable only by a party holding the shared symmetric key, and FTISCON's
smart contract check still trusts the controller that requested the
installation in the first place, so a controller that is itself compromised
and issues a policy consistent but malicious rule evades both.

Detection frameworks that bring mobility awareness or collaborative training
to this space still fall short of covering hidden forwarding specifically.
Liu et al. propose a blockchain and federated learning framework for
collaborative intrusion detection that trains across distributed vehicular
edge nodes rather than a central server, which is directly relevant to
preserving the location privacy a hidden forwarding defense must itself not
undermine while collecting evidence \cite{Liu2021BFLVEC}. Post quantum
authentication protocols for the vehicle to infrastructure exchange, such as
the lattice based scheme of El Dalahmeh and Alia, secure the credential
exchange but detect nothing about what happens to a packet after it is
accepted onto the network \cite{ElDalahmeh2025}. Neither line of work
addresses the passive case at all, where the duplicated packet is a genuinely
authentic retransmission that no cryptographic check, however strong, can
distinguish from a legitimate second delivery.

Table~\ref{tab:vanguard_related} compares six representative works against
five properties that determine suitability for this paper's setting: whether
the design detects passive, cryptographically indistinguishable duplication
specifically, rather than only active, forged forwarding, whether it uses
post quantum cryptography, whether detection is federated, whether it
preserves location privacy by design rather than as an afterthought, and what
trust it places in the controller.

\begin{table*}[!t]
	\centering
	\caption{Comparison of representative defenses relevant to hidden forwarding attacks in software defined vehicular networks.}
	\label{tab:vanguard_related}
	\renewcommand{\arraystretch}{1.25}
	\footnotesize
	\begin{tabular}{|>{\raggedright\arraybackslash}p{2.3cm}|>{\centering\arraybackslash}p{2.0cm}|>{\centering\arraybackslash}p{1.6cm}|>{\centering\arraybackslash}p{1.6cm}|>{\centering\arraybackslash}p{1.8cm}|>{\raggedright\arraybackslash}p{2.4cm}|}
		\hline
		\textbf{Work} & \makecell{\textbf{Detects}\\\textbf{Passive Copy}} & \makecell{\textbf{Post}\\\textbf{Quantum}} & \textbf{Federated} & \makecell{\textbf{Privacy}\\\textbf{Preserving}} & \textbf{Controller Trust} \\
		\hline
		HSA~\cite{Mohammadi2016} & No & No & No & No & Assumed trusted \\
		\hline
		FADE~\cite{Pang2016FADE} & No & No & No & No & Assumed trusted \\
		\hline
		SDNsec~\cite{Sasaki2016SDNsec} & No & No & No & No & Assumed trusted \\
		\hline
		FTISCON~\cite{Bargayary2023FTISCON} & No & No & No & No & Blockchain verified \\
		\hline
		Liu et al.~\cite{Liu2021BFLVEC} & No & No & Yes & Partial & Assumed trusted \\
		\hline
		Post quantum V2R authentication~\cite{ElDalahmeh2025} & Not applicable & Yes & No & Yes & Assumed trusted \\
		\hline
		\textbf{VANGUARD HF (proposed)} & \textbf{Yes} & \textbf{Yes} & \textbf{Yes} & \textbf{Yes} & \textbf{Zero trust, controller included} \\
		\hline
	\end{tabular}
\end{table*}

A related, learning based line of work targets a compromised switch's
behaviour directly rather than a specific forwarding signature. Huang et
al. propose FSDM, extracting the distribution of control channel
occupation rate to detect control plane saturation \cite{Huang2020}, and
a later enhancement frames the same problem as multivariate time series
anomaly detection over a compromised switch's general behaviour,
including delaying, dropping, or modifying traffic \cite{Chi2020}. Both
are saturation or general misbehaviour detectors rather than duplication
specific, and neither reads a signal comparable to the reception log
anomaly score this paper introduces. Federated approaches extended toward
zero day coverage report encouraging results on adjacent problems: an
edge based framework combining federated learning with blockchain reports
high accuracy against previously unseen attack types in intelligent
transportation systems \cite{AbouElHouda2024}, and a comparison of
recurrent architectures under federated training on the VeReMi dataset
reports precision between 68 and 94 percent across attack types
\cite{Mansouri2024}, though neither targets duplication or extends
distrust to the controller. The selective time delay literature, slow
flow table exhaustion \cite{Pascoal2017,Pascoal2020DoS,Cao2022LOFT}, a
slow rate specific detector \cite{Tang2023SFTOGuard}, protocol level
timing prevention \cite{Arsalan2018}, and mobility aware misbehavior
trust scoring \cite{Nayak2022}, addresses a disjoint attack class from
the one this paper considers and is reviewed in full in the companion
paper PHANTOM \cite{Phantom2026}.

The passive copy column in Table~\ref{tab:vanguard_related} is empty for
every prior work. Path validation and flow conservation approaches test
observable forwarding behaviour, which a passive duplicate does not disturb,
since the original packet is delivered exactly as expected while a
retransmitted, genuinely authentic copy departs unnoticed. Cryptographic
accountability schemes verify that a packet was forwarded correctly, which a
passive duplicate also satisfies, since nothing about the packet itself was
altered. This is the gap VANGUARD HF addresses: detection that does not rely
on any observable disturbance to the legitimate path at all, achieved through
independent witness corroboration from nearby vehicles rather than through
inspecting the packet or the flow it belongs to, under a threat model that
extends zero trust to the controller alongside every roadside unit and
vehicle. The companion paper PHANTOM addresses the disjoint attack class this
project also considers, selective time delay, under the same architecture
\cite{Phantom2026}; the joint evaluation of both classes active
simultaneously is reported separately in NEXUS \cite{Nexus2026}.


\section{System Model and Attack Formalization}
\label{system_model}

\subsection{Network and Threat Model}
\label{sec:threat_model}

The network comprises a mobile vehicle population, a fixed grid of
roadside units, and a flat, registered set of controllers, connected as
formalised in Section~\ref{sec:ctrl_trust}. Controllers issue flow
modification instructions that roadside units install and enforce, and
vehicles communicate over both vehicle to infrastructure and vehicle to
vehicle paths, either of which a hidden forwarding attack can exploit to
place a duplicate.

The threat model assumes an internal or external adversary capable of
delaying, duplicating, fabricating, or covertly forwarding packets by
exploiting flow modification installation and forwarding behaviour, from
either the control plane or the data plane, never both at once from a
single compromised entity. The attacker favours low rate, low visibility
techniques that avoid a volumetric alarm, which for this paper's own
attack class means the original packet is left completely undisturbed on
its legitimate path, since the moment delivery itself looks abnormal the
element of covert observation that makes hidden forwarding worthwhile is
already lost. Zero trust extends to the controller itself, not only to
roadside units and vehicles: every controller in the registered set is
granted read only blockchain access, and its behaviour is audited
exclusively through independent roadside unit logging rather than any
form of self reporting. Roadside units are individually corruptible but
collectively honest under a Byzantine fault tolerance bound of $f<n/3$,
and no controller is permanently active or inherently trusted regardless
of how long it has been in service. This one bound $f$ is the parameter
behind every threshold expressed as $f{+}1$ or $2f{+}1$ later in this
paper, whether the quorum being reached is roadside unit endorsement or
independent witness confirmation, rather than a separate value chosen
independently for each mechanism.

\subsection{Attack Scenarios}
\label{sec:attack_scenarios}

Hidden forwarding leaves the original packet on its legitimate path
entirely undisturbed while a duplicate is secretly routed to an
unauthorized destination, enabling eavesdropping without any observable
disruption to normal delivery. Figure~\ref{fig:hf_attacks} shows all four
variants this paper addresses, split by which plane the adversary
occupies and by whether the duplicate is actively fabricated with a
forged signature or passively retransmitted carrying the original,
genuinely valid one.
\begin{figure}[!t]
	\centering
	\includegraphics[width=0.95\linewidth]{Hidden Forwarding Attack Versions1.png}
	\caption{Hidden Forwarding Attack Variants}
	\label{fig:hf_attacks}
\end{figure}
In Signature S5, the active control plane variant, a compromised
controller issues a flow modification that installs an unauthorized
routing path at the roadside unit; once installed, the roadside unit
forwards a targeted packet to its correct, original destination exactly
as expected while simultaneously fabricating and sending a duplicate to
an unauthorized listener, so nothing about the legitimate delivery ever
looks wrong. In Signature S6, the active data plane variant, the
roadside unit or an attacking vehicle is itself compromised and performs
the same fabrication directly, without any corresponding controller
instruction. Signatures S7 and S8 achieve the passive form of the same
outcome, a controller, for S7, or a compromised roadside unit or vehicle,
for S8, retransmits a genuinely authentic copy of the original packet
itself rather than fabricating one, so the duplicate carries a valid
signature over unmodified content and every cryptographic check available
to a verifier passes, since nothing about the packet was actually forged.
Table~\ref{tab:vanguard_taxonomy} summarises all four variants against
the plane each occupies, the property that makes each one covert, and
the primary mechanism this paper assigns to detect it.
\begin{table}[!t]
	\centering
	\caption{Hidden Forwarding Attack Taxonomy}
	\label{tab:vanguard_taxonomy}
	\renewcommand{\arraystretch}{1.2}
	\footnotesize
	\begin{tabular}{|c|c|>{\raggedright\arraybackslash}p{2.8cm}|>{\raggedright\arraybackslash}p{1.9cm}|}
		\hline
		\textbf{Variant} & \textbf{Plane} & \textbf{Covertness} & \textbf{Detector} \\
		\hline
		S5 & Control & Original path unaffected, duplicate fabricated & Signature S5 \\
		\hline
		S6 & Data & Original path unaffected, duplicate fabricated & Signature S6 \\
		\hline
		S7 & Control & Duplicate is genuinely authentic, passes all checks & Classifier, witness quorum \\
		\hline
		S8 & Data & Duplicate is genuinely authentic, passes all checks & Classifier, witness quorum \\
		\hline
	\end{tabular}
\end{table}
Each signature's formal detection condition, corrected to require that
the accusing node was genuinely addressed the duplicate rather than a
bystander that merely overheard it, is developed in
Section~\ref{sec:dual_mode}, where the detection engine that evaluates
it online is presented; this section formalises what each attack does,
not yet how it is caught. Signatures S1 through S4, selective time
delay, are a disjoint attack class developed independently in the
companion paper PHANTOM \cite{Phantom2026}.

\subsection{Mobility Induced Threat Amplification}
\label{sec:mobility_amplification}

Hidden forwarding is masked by vehicular mobility for a reason distinct
from the timing based masking the companion paper addresses. Flow
conservation detectors such as HSA and FADE, reviewed in
Section~\ref{sec:related}, count packets over a flow's observed lifetime
to test whether the volume entering and leaving remains consistent, and
that test needs enough samples to be statistically reliable. A vehicle's
residence time inside a single roadside unit's coverage zone bounds how
many samples such a detector ever gets to see, given in
Equation~\eqref{eq:observation_window}.
\begin{equation}
	N_{obs} = r \cdot t_{zone} < N_{min}(\alpha)
	\label{eq:observation_window}
\end{equation}
As given in Equation~\eqref{eq:observation_window}, $N_{obs}$ is the
number of packets observable during a single zone crossing, the product
of the per flow packet rate $r$ and the zone residence time $t_{zone}$,
typically nine to twenty two seconds at highway speed. Once this falls
below the minimum sample size $N_{min}(\alpha)$ a flow conservation test
needs at significance level $\alpha$, such a detector becomes
statistically unreliable regardless of how its threshold is tuned,
independent of how the duplication itself is carried out. This is what
motivates per packet, per hop cryptographic verification over flow level
statistics as the primary mechanism developed in
Section~\ref{sec:crypto_layer}, since a zero knowledge proof requires no
observation window at all, complemented by the federated classifier of
Section~\ref{sec:fed_lstm} and the witness quorum of
Section~\ref{sec:witness} for the passive case a proof alone cannot
cover. The three mechanisms amplifying selective time delay specifically,
handoff masking, topology staleness, and flow rule churn camouflage, do
not apply to the attack class this paper addresses and are developed in
the companion paper PHANTOM \cite{Phantom2026}.


\section{Proposed Framework}
\label{sec:framework}

\subsection{Multi Controller Zero Trust Architecture}
\label{sec:ctrl_trust}

VANGUARD HF shares its controller architecture with the companion paper
PHANTOM, since a compromised controller can launch either attack class from
the same position of trust, and a defense that only distrusts the controller
for one class is not a zero trust defense at all \cite{Phantom2026}. A flat,
registered controller set $\mathcal{C} = \{c_1, \ldots, c_m\}$ replaces any
single, permanently trusted controller, with the trusted subset at time $t$
given in Equation~\eqref{eq:trusted_ctrl_set}.
\begin{equation}
	\mathcal{C}_{trusted}(t) = \{c_i \in \mathcal{C} : T_{c_i}(t) \geq T_{min}^{ctrl}\}
	\label{eq:trusted_ctrl_set}
\end{equation}
As given in Equation~\eqref{eq:trusted_ctrl_set}, a controller whose trust
falls below the operational minimum is excluded from serving roadside units
regardless of how it came to fall below that threshold. Assignment of
roadside units to a trusted controller is given in
Equation~\eqref{eq:rsu_ctrl_assign}.
\begin{equation}
	c^*(k,t) = \underset{c_i \in \mathcal{C}_{trusted}(t)}{\arg\min}\, d(r_k, c_i)
	\label{eq:rsu_ctrl_assign}
\end{equation}
As given in Equation~\eqref{eq:rsu_ctrl_assign}, $d(r_k,c_i)$ is the
geographic distance between roadside unit $r_k$ and controller $c_i$, so
the assignment is recomputed over whichever controllers are currently
trusted rather than fixed once at deployment time. Controller trust is penalised on two channels. Conflict evidence
accumulates whenever a roadside unit receives a flow modification absent
from the blockchain committed policy of Section~\ref{sec:blockchain_arch},
which for this paper is the same check Signature S5 also relies on. A
second, independent channel keyed to Signature S1 additionally penalises a
controller for timing manipulation even when every flow modification it
issues is individually authorized; this channel is PHANTOM's own
contribution and is developed there in full, since it is keyed to a
signature this paper does not otherwise use \cite{Phantom2026}. The
resulting update rule, shared by both evidence channels, is given in
Equation~\eqref{eq:ctrl_trust_update}.
\begin{equation}
	T_{c_i}^{(t+1)} =
	\begin{cases}
		\min(1, T_{c_i}^{(t)} + \Delta_r^{ctrl}) & \text{conflict evidence} < f{+}1 \;\land\; \mathcal{E}_{delay}(c_i,t) < f{+}1 \\
		\max(0, T_{c_i}^{(t)} - \Delta_p^{ctrl}) & \text{conflict evidence} \geq f{+}1 \;\lor\; \mathcal{E}_{delay}(c_i,t) \geq f{+}1
	\end{cases}
	\label{eq:ctrl_trust_update}
\end{equation}
As given in Equation~\eqref{eq:ctrl_trust_update}, a threshold of $f{+}1$
independent roadside units is required on either channel before a penalty
applies, so a single compromised roadside unit cannot unilaterally revoke a
controller through either evidence path. Revocation and failover follow
automatically once trust drops below the operational threshold, given in
Equations~\eqref{eq:sc_revoke} and~\eqref{eq:ctrl_failover}.
\begin{equation}
	\textsc{SC.Revoke}(c_i) \Leftarrow T_{c_i}^{(t)} < T_{min}^{ctrl}
	\label{eq:sc_revoke}
\end{equation}
\begin{equation}
	c^*(k, t^+) = \underset{c_i \in \mathcal{C}_{trusted}(t^+)}{\arg\min}\, d(r_k, c_i), \quad \mathcal{C}_{trusted}(t^+) = \mathcal{C}_{trusted}(t) \setminus \{c_i^{revoked}\}
	\label{eq:ctrl_failover}
\end{equation}
As given in Equations~\eqref{eq:sc_revoke} and~\eqref{eq:ctrl_failover}, no
pre designated backup controller is needed, since every affected roadside
unit re executes the assignment over whatever trusted controllers remain.

The zero knowledge proof key material this paper's own hop legitimacy proof
depends on is generated independently of the controller through a
distributed key generation ceremony among roadside unit ring members, with
the resulting verification key committed as given in
Equation~\eqref{eq:vk_commit}.
\begin{equation}
	\textsc{BC.Commit}\!\left(vk_{ZKP}, \{\text{Com}_j\}_{j=1}^{n_{RSU}}, ts_{setup}\right)
	\label{eq:vk_commit}
\end{equation}
As given in Equation~\eqref{eq:vk_commit}, $vk_{ZKP}$ is the jointly
generated verification key the hop legitimacy proof of
Section~\ref{sec:crypto_layer} is checked against, and $\{\text{Com}_j\}$
are the per roadside unit commitments from the ceremony, so a compromised
controller has no path to influencing the key plane this proof depends on.
When a participating roadside unit is later revoked, every proving key it
contributed to is rotated rather than left in circulation, triggered as
given in Equation~\eqref{eq:key_rotation_trigger} and recommitted as given
in Equation~\eqref{eq:vk_commit_rotated}.
\begin{equation}
	\textsc{SC.RotateKeys} \Leftarrow \exists\, r_j : T_{r_j}^{(t)} < T_{min} \land r_j \in \mathcal{P}_{DKG}^{current}
	\label{eq:key_rotation_trigger}
\end{equation}
\begin{equation}
	\textsc{BC.Commit}\!\left(vk_{ZKP}^{new}, \{\text{Com}_j\}_{j \in \mathcal{P}_{DKG}^{new}}, ts_{rotation}\right)
	\label{eq:vk_commit_rotated}
\end{equation}
As given in Equation~\eqref{eq:vk_commit_rotated}, $\mathcal{P}_{DKG}^{new}$
is the prior participant set with the revoked roadside unit removed, so
rotation never requires a full ceremony among every roadside unit in the
network, a property equally important here as in PHANTOM since the hop
legitimacy proof this paper relies on uses the same key plane.


\subsection{Dual Mode Detection Engine}
\label{sec:dual_mode}

Figure~\ref{fig:dual_mode} shows the two cooperating detection paths this
paper shares with PHANTOM; this section develops the roadside unit path in
full, since Signatures S5 through S8 are evaluated there exclusively,
following escalation from whichever on board signatures fired.
\begin{figure}[!t]
	\centering
	\includegraphics[width=0.95\linewidth]{Figure_1_Final.eps}
	\caption{Dual Mode Framework Overview}
	\label{fig:dual_mode}
\end{figure}

Hidden forwarding is masked by vehicular mobility for the reason
Section~\ref{sec:mobility_amplification} establishes: a vehicle's
residence time inside a single roadside unit's zone bounds how many
packets a flow conservation detector ever gets to observe, given in
Equation~\eqref{eq:observation_window}, which is what motivates per
packet cryptographic verification over flow level statistics as the
primary mechanism developed in this section, since a zero knowledge
proof requires no observation window at all.

Active hidden forwarding fabricates a duplicate that a verifier can reject
cryptographically; passive hidden forwarding retransmits the original,
genuinely authentic packet, so no cryptographic check can reject it and
detection instead depends on an addressing property read from the link
layer rather than from content. Signature S5, the active control plane
variant, is given in Equation~\eqref{eq:sig_s5}.
\begin{equation}
	\begin{aligned}
		\text{S5:}\quad
		& d' \notin \mathcal{P}(s,d) \\
		& \wedge\; \neg\,\textsc{BC.Endorsed}\big(\texttt{FlowMod}(s \to d')\big) \\
		& \wedge\; \textsc{ML-DSA-87.Verify}(\sigma_{copy}, pk_s, m_{copy}) = 0 \\
		& \wedge\; b_{hop}(u) = 0 \\
		& \wedge\; \textsc{Addressed}(p, d') = 1
	\end{aligned}
	\label{eq:sig_s5}
\end{equation}
As given in Equation~\eqref{eq:sig_s5}, the second conjunct establishes the
control plane origin through the absence of an endorsed policy for the flow
toward $d'$ rather than through aggregate signature failure: a fabricated
duplicate delivered off path is never part of the primary path signing
batch in the first place, so it cannot fail a check it was never subject to,
and requiring it would make the signature unsatisfiable. The addressing
conjunct is read directly from the link layer destination of the frame the
accusing node received, independent of the signed next hop, and is what
separates a node that was genuinely sent the unauthorized duplicate from a
bystander that merely overheard a broadcast transmission on the shared
radio medium, since every neighbour within range receives every frame.
Signature S6, the active data plane variant, is given in
Equation~\eqref{eq:sig_s6}.
\begin{equation}
	\begin{aligned}
		\text{S6:}\quad
		& \exists\, d, d' : d \neq d' \\
		& \wedge\; H(p) \in \mathcal{R}(d, W_{S6}) \\
		& \wedge\; \textsc{BC.QueryReceipt}(H(p), d', W_{S6}) = 1 \\
		& \wedge\; \textsc{ML-DSA-87.Verify}(\sigma_{copy}, pk_s, m_{copy}) = 0 \\
		& \wedge\; b_{hop}(u) = 0 \\
		& \wedge\; \textsc{Addressed}(p, d') = 1
	\end{aligned}
	\label{eq:sig_s6}
\end{equation}
As given in Equation~\eqref{eq:sig_s6}, observing the same packet hash $H(p)$
locally at $d$ and, via the blockchain receipt query of
Equation~\eqref{eq:packet_receipt_log}, at $d'$ within window $W_{S6}{=}30$s
confirms duplication; the window must exceed the ten second evaluation block
this paper's own evaluation methodology uses, since a shorter window can
observe the two receptions in adjacent, disconnected blocks, though a window
this long frequently spans a roadside unit handoff, a limitation of the
signature rather than of the window choice itself. Signature S6 carries no
flow modification absence conjunct: an earlier construction included one to
separate S6 from S5, but the underlying endorsement query reported the
monitored flow as endorsed in every cycle, which suppressed every correct S6
detection rather than only S5 events, and it is not part of the signature
for that reason. Signature S7, the passive control plane variant, is given
in Equation~\eqref{eq:sig_s7}.
\begin{equation}
	\text{S7:}\; \tfrac{d}{dt}\text{Vol}(d',t) > \epsilon_{vol} \;\land\; \textsc{ML-DSA-87.Verify}(\sigma_c, pk_s, m_c) = 1 \;\land\; b_{hop}(u) = 0 \;\land\; \textsc{Addressed}(p, d') = 1
	\label{eq:sig_s7}
\end{equation}
As given in Equation~\eqref{eq:sig_s7}, verification succeeding on the copy
is what distinguishes this passive variant from S5's active fabrication: the
content is unmodified because the original, genuinely signed packet was
retransmitted rather than forged, so cryptographic evidence alone cannot be
the primary signal here, and the witness mechanism of
Section~\ref{sec:witness} carries that role instead, with the volume rate
and hop proof failure serving only as corroborating indicators. The volume
rate is estimated over a tumbling ten second window with
$\epsilon_{vol}{=}0.1$ packets per second. Signature S8, the passive data
plane variant, is given in Equation~\eqref{eq:sig_s8}.
\begin{equation}
	\text{S8:}\; \textsc{BatchVerify}(\boldsymbol{\sigma}, \{pk_i\}, \{m_i\}, \mathbf{r}) = 1 \;\land\; \textsc{ML-DSA-87.Verify}(\sigma_c, pk_s, m_c) = 1 \;\land\; b_{hop}(u) = 0 \;\land\; \textsc{Addressed}(p, d') = 1
	\label{eq:sig_s8}
\end{equation}
As given in Equation~\eqref{eq:sig_s8}, aggregate verification succeeding
confirms correct delivery across every hop on the primary path, so, exactly
as for S7, the primary detection mechanism is witness monitoring rather than
this signature alone.

Algorithm~\ref{alg:lrad_rsu} summarises the roadside unit procedure. It is
presented here corrected against two inconsistencies confirmed directly in
the current draft: the addressing conjunct now required by
Equations~\eqref{eq:sig_s5} through~\eqref{eq:sig_s8} is included in every
relevant step, the stale flow modification absence check is removed from
the S6 and S7 steps to match the equations above, the composite decision is
stated as confirmed against the live implementation rather than the draft
text's own, differently structured version, and the reception log anomaly
rule verified in Section~\ref{sec:witness} is included as a direct
contributor to Signatures S7 and S8 rather than only as a condition
suppressing the classifier.
\begin{algorithm}[!t]
	\caption{Lightweight and Full Rule Based Attack Detection at the Roadside Unit (LRAD RSU), Hidden Forwarding Steps}
	\label{alg:lrad_rsu}
	\begin{algorithmic}[1]
		\small
		\Procedure{LRAD-RSU (S5--S8 excerpt)}{$p, r, t, \boldsymbol{\sigma}, \pi_{delay}, \pi_{hop}$}
		\State $\mathbf{r} \gets H_{\text{SHA3-512}}(\boldsymbol{\sigma}\|m_1\|\cdots\|m_n)$ \Comment{Eq.~\eqref{eq:batch_challenge}}
		\State $b_{batch} \gets \textsc{BatchVerify}(\boldsymbol{\sigma}, \{pk_i\}, \{m_i\}, \mathbf{r})$ \Comment{Eq.~\eqref{eq:batch_verify}}
		\State $b_{hop} \gets \textsc{STARK.Verify}(\pi_{hop}, C_{hop}, H_{\text{SHA3-512}})$ \Comment{Eq.~\eqref{eq:stark_hop_verify}}
		\If{$\neg b_{batch}$}
		\State $v_{atk,copy} \gets \arg\min_{i}\{i : \textsc{ML-DSA-87.Verify}(\sigma_i,pk_i,m_i)=0\}$ \Comment{Eq.~\eqref{eq:batch_fallback}}
		\EndIf
		\State $\text{CopyVerify}_{d'} \gets \textsc{BC.QueryReceipt}(H(p),d',W).\text{verify}$ \Comment{Eq.~\eqref{eq:packet_receipt_log}}
		\State $flag_{S5} \gets [\neg b_{batch}] \land [f_{unauth}(r,t){=}1] \land [\text{CopyVerify}_{d'}{=}0] \land [b_{hop}{=}0] \land [\textsc{Addressed}(p,d'){=}1]$ \Comment{Eq.~\eqref{eq:sig_s5}}
		\State $\text{RecvAt}_{d'} \gets \textsc{BC.QueryReceipt}(H(p),d',W_{S6})$
		\State $flag_{S6} \gets [\text{RecvAt}_{d'}{=}1] \land [\text{CopyVerify}_{d'}{=}0] \land [b_{hop}{=}0] \land [\textsc{Addressed}(p,d'){=}1]$ \Comment{Eq.~\eqref{eq:sig_s6}, no FlowMod absence term}
		\State $flag_{S7} \gets [\tfrac{d}{dt}\text{Vol}(d',t){>}\epsilon_{vol}] \land [\text{CopyVerify}_{d'}{=}1] \land [b_{hop}{=}0] \land [\textsc{Addressed}(p,d'){=}1]$ \Comment{Eq.~\eqref{eq:sig_s7}, no FlowMod absence term}
		\State $flag_{S8} \gets [b_{batch}] \land [\text{CopyVerify}_{d'}{=}1] \land [b_{hop}{=}0] \land [\textsc{Addressed}(p,d'){=}1]$ \Comment{Eq.~\eqref{eq:sig_s8}}
		\State $flag_{S7} \gets flag_{S7} \lor [R_{anom}(r,t) > 0]$; \quad $flag_{S8} \gets flag_{S8} \lor [R_{anom}(r,t) > 0]$ \Comment{Eq.~\eqref{eq:feat_ranom}, direct contribution}
		\If{$\neg(flag_{S3} \lor flag_{S4} \lor [R_{anom}(r,t){>}0])$}
		\State $flag_{LSTM} \gets D_{LSTM}^{(k)}$ \Comment{Eq.~\eqref{eq:lstm_gate}}
		\Else
		\State $flag_{LSTM} \gets 0$
		\EndIf
		\State $D_{RSU} \gets flag_{S5} \lor flag_{S6} \lor flag_{S7} \lor flag_{S8} \lor flag_{LSTM}$
		\If{$D_{RSU}$ and once contained, $flag_{S5}, \ldots, flag_{S8} \gets 0$ for that node} \Comment{Quarantine suppression, Section~\ref{sec:witness}}
		\EndIf
		\If{$D_{RSU}$}
		\State \texttt{BTMM}$(\text{attributed node}, \text{plane})$ \Comment{Section~\ref{sec:blockchain_arch}}
		\EndIf
		\State \textbf{return} $D_{RSU}$
		\EndProcedure
	\end{algorithmic}
\end{algorithm}
As shown in Algorithm~\ref{alg:lrad_rsu}, batch and hop verification are
computed once per packet and reused across all four signatures rather than
recomputed per signature, and the reception log anomaly rule contributes
directly to the two passive signatures rather than only through the
classifier path, closing a gap that existed while this rule ran only as a
classifier suppression condition. Signatures S1 through S4 are evaluated in
the same procedure and are developed in full in PHANTOM
\cite{Phantom2026}.


\subsection{Post Quantum Cryptographic Layer}
\label{sec:crypto_layer}

Figure~\ref{fig:crypto_pipeline} shows the cryptographic pipeline every
packet carries; this paper develops the hop legitimacy proof in full, since
it is the primary cryptographic contributor to hidden forwarding detection.
\begin{figure}[!t]
	\centering
	\includegraphics[width=0.95\linewidth]{Figure_2_Final.eps}
	\caption{Cryptographic Packet Pipeline}
	\label{fig:crypto_pipeline}
\end{figure}

Every forwarding node signs each packet it forwards, given in
Equation~\eqref{eq:mldsa_sign}.
\begin{equation}
	\sigma_i = \textsc{ML-DSA-87.Sign}\!\left(sk_i,\; H_{\text{SHA3-512}}\!\left(\text{msg\_id} \,\|\, ts_i \,\|\, \eta_i \,\|\, \text{nh}_i \,\|\, z_i\right)\right)
	\label{eq:mldsa_sign}
\end{equation}
As given in Equation~\eqref{eq:mldsa_sign}, the signed message binds the
packet identifier, the forwarding timestamp $ts_i$, a fresh nonce
$\eta_i$, the node's chosen next hop $\text{nh}_i$, and its current zone
$z_i$ under the node's own key $sk_i$, so a signature over one next hop
cannot later be replayed to vouch for a different one, which is exactly
the property Signature S5's forged copy has no way to produce honestly.
Resource constrained on board units instead use a symmetric key tag, given
in Equation~\eqref{eq:hmac_light}.
\begin{equation}
	\tau_i = \textsc{HMAC-SHA3-512}\!\left(k_i,\; \text{msg\_id} \,\|\, ts_i \,\|\, \eta_i\right)
	\label{eq:hmac_light}
\end{equation}
As given in Equation~\eqref{eq:hmac_light}, $\tau_i$ omits the next hop
and zone fields entirely, since a lightweight node's own tag is only ever
checked by the roadside unit holding the matching key $k_i$, not
propagated as evidence to a third party the way a full signature is.
Verifying every hop's signature independently does not scale to a multi hop
path, so a randomized challenge permits one combined check instead, given in
Equations~\eqref{eq:batch_challenge} and~\eqref{eq:batch_verify}.
\begin{equation}
	\mathbf{r} = (r_1, \ldots, r_n) = H_{\text{SHA3-512}}\!\left(\sigma_1 \,\|\, \cdots \,\|\, \sigma_n \,\|\, m_1 \,\|\, \cdots \,\|\, m_n\right)
	\label{eq:batch_challenge}
\end{equation}
As given in Equation~\eqref{eq:batch_challenge}, the challenge vector
$\mathbf{r}$ is derived only after every signature and message on the
path has already been produced, so an adversary cannot choose a forged
signature to favourably match a challenge it has not yet seen.
\begin{equation}
	\textsc{BatchVerify}\!\left(\{\sigma_i\}, \{pk_i\}, \{m_i\}, \mathbf{r}\right) = \left[\sum_{i=1}^{n} r_i \mathbf{A}_i \mathbf{z}_i \stackrel{?}{=} \sum_{i=1}^{n} r_i \left(\mathbf{t}_i \mathbf{c}_i + \mathbf{w}'_i\right) \pmod{q}\right]
	\label{eq:batch_verify}
\end{equation}
As given in Equation~\eqref{eq:batch_verify}, one combined lattice equation
accepts or rejects all $n$ signatures at once, which is what Signature S8
tests directly and what Signatures S5 and S6 test the failure of. When the
batch check fails, the specific forger is isolated as given in
Equation~\eqref{eq:batch_fallback}.
\begin{equation}
	v_{atk} = \arg\min_{i \in \{1,\ldots,n\}} \Bigl\{i : \textsc{ML-DSA-87.Verify}(\sigma_i, pk_i, m_i) = 0\Bigr\}
	\label{eq:batch_fallback}
\end{equation}
As given in Equation~\eqref{eq:batch_fallback}, individual verification
runs only on the batch that already failed, so the cost of finding $v_{atk}$
is paid exactly once per genuine forgery rather than on every packet that
would have passed the combined check anyway.

Next hop legitimacy is certified by a zero knowledge proof rather than
trusting a forwarding node's own account of its routing choice. The proof
is given in Equation~\eqref{eq:stark_hop}.
\begin{equation}
	\pi_{hop,i} = \textsc{STARK.Prove}\!\left(\{\sigma^{rcpt}_{i+1},\, \mathcal{P}_i\};\; C_{hop};\; H_{\text{SHA3-512}}\right)
	\label{eq:stark_hop}
\end{equation}
As given in Equation~\eqref{eq:stark_hop}, the private witnesses are the
receipt signature from the next hop and the local copy of the authorized
next hop policy set committed at flow setup; the arithmetic circuit $C_{hop}$
encodes that the receipt signature verifies under the claimed next hop's key
and that the claimed next hop lies within the authorized policy set,
certifying both delivery and routing compliance without revealing the
receipt signature itself to any external verifier. Verification is given in
Equation~\eqref{eq:stark_hop_verify}.
\begin{equation}
	b_{hop,i} = \textsc{STARK.Verify}\!\left(\pi_{hop,i},\; C_{hop},\; H_{\text{SHA3-512}}\right) \in \{0, 1\}
	\label{eq:stark_hop_verify}
\end{equation}
As given in Equation~\eqref{eq:stark_hop_verify}, $b_{hop,i}{=}0$ is a
corroborating cryptographic indicator for both active and passive hidden
forwarding, never the primary signal, since a compromised node controls its
own proof generation and can produce a valid proof while covertly
duplicating over the shared radio medium; this is exactly why the witness
mechanism of Section~\ref{sec:witness} and the classifier of
Section~\ref{sec:fed_lstm} exist as independent detection paths rather than
relying on this proof alone. The companion timing compliance proof,
$\pi_{delay}$, is generated at the same pipeline stage and is developed in
full in PHANTOM, since it is the primary cryptographic signal there
\cite{Phantom2026}.


\subsection{Federated LSTM with Byzantine Robust Aggregation}
\label{sec:fed_lstm}

Rule based and cryptographic signatures cover every variant with a
deterministic or near deterministic observable, but passive hidden
forwarding leaves none: the retransmitted packet is a genuinely authentic
copy, so every cryptographic check available passes. A recurrent classifier
trained across roadside units without centralizing raw traffic provides the
complementary, learned detection path this residual case requires.
Figure~\ref{fig:fed_lstm} shows the architecture together with the
aggregation procedure described below.
\begin{figure}[!t]
	\centering
	\includegraphics[width=0.95\linewidth]{Figure_3_Final.eps}
	\caption{Federated LSTM Architecture}
	\label{fig:fed_lstm}
\end{figure}

Each roadside unit observes an eleven feature vector per cycle, given in
Equation~\eqref{eq:lstm_input}.
\begin{equation}
	\mathbf{x}_t^{(r)} = \left[
	\delta^{\max}_t,\;
	\lambda_{PI,t},\;
	U_{TCAM,t},\;
	\mathbb{1}[\pi_{delay}{=}\bot],\;
	\mathbb{1}[\pi_{hop}{=}\bot],\;
	\rho_t,\;
	\bar{v}_t,\;
	\log(1{+}N_{div,t}),\;
	\log(1{+}A^{\mathrm{RSU}}_{tp,t}),\;
	\log(1{+}R_{anom,t}),\;
	\mathbb{1}[\delta^{\max}_t{>}\Delta_{max}]
	\right]^T
	\label{eq:lstm_input}
\end{equation}
As given in Equation~\eqref{eq:lstm_input}, three of the eleven features are
specific to hidden forwarding and are developed in full below; each is
transformed by $\log(1{+}x)$ before entering the vector, since its benign
resting value is small and constant while its value under attack spans
several orders of magnitude. The per roadside unit destination diversity
count is given in Equation~\eqref{eq:feat_ddiv}.
\begin{equation}
	N_{div}(r,t) = \left|\left\{\, v \in \mathcal{V}(r,t) \;:\; \mathcal{D}(v,r,t) \not\subseteq \mathcal{P}(v,\cdot) \,\right\}\right|
	\label{eq:feat_ddiv}
\end{equation}
As given in Equation~\eqref{eq:feat_ddiv}, $N_{div}(r,t)$ counts the distinct
source vehicles, among those whose deliveries attribute to roadside unit
$r$, whose reached destinations include at least one outside their
authorized policy, with each delivery attributed to the covering roadside
unit of its last hop sender. Unlike the reception log anomaly score below,
this count separates one high volume attacker from several quiet ones,
correlating with it at only 0.47 to 0.69 across variants, so the two carry
genuinely different information rather than one being redundant with the
other. When a roadside unit relays no delivery of the monitored flow in a
given cycle, the feature takes the resting value one rather than the count
itself. The per flow throughput asymmetry is given in
Equation~\eqref{eq:feat_atp}.
\begin{equation}
	A_{tp}(v,r,t) = \frac{\lambda_{bytes}(v,d,W)}{\lambda_{bytes}(v,\cdot,W)}
	\label{eq:feat_atp}
\end{equation}
As given in Equation~\eqref{eq:feat_atp}, this ratio falls below one under
active hidden forwarding, since traffic is diverted, and rises above one
under passive hidden forwarding, since attributed throughput now includes
the retransmitted copy. Since the detector consumes one vector per roadside
unit per cycle, this per source ratio is supplied as the per roadside unit
aggregate given in Equation~\eqref{eq:feat_atp_rsu}.
\begin{equation}
	A_{tp}^{\mathrm{RSU}}(r,t) = \frac{\left|\left\{p \in \mathcal{F}_m(t) : c(u_p,t)=r,\; dst(p)=d_p\right\}\right|}{\left|\left\{p \in \mathcal{F}_m(t) : c(u_p,t)=r\right\}\right|}
	\label{eq:feat_atp_rsu}
\end{equation}
As given in Equation~\eqref{eq:feat_atp_rsu}, $\mathcal{F}_m(t)$ is the set
of monitored flow deliveries in the cycle, $u_p$ the last hop sender of
delivery $p$, and $c(u,t)$ the roadside unit currently covering $u$; each
authorized delivery is credited to that covering roadside unit rather than
computed as one value shared across the network, correcting a broadcast
scalar implementation that previously carried zero per roadside unit
information. The reception log anomaly score, the most novel of the three,
is given in Equation~\eqref{eq:feat_ranom}.
\begin{equation}
	R_{anom}(r,t) = \frac{1}{W}\left|\left\{ H(p) : \textsc{BC.QueryReceipt}(H(p),d',W) = 1,\; d' \notin \mathcal{P}(s_p, d_p) \right\}\right|
	\label{eq:feat_ranom}
\end{equation}
As given in Equation~\eqref{eq:feat_ranom}, this counts distinct packet
hashes per unit time received at unauthorized destinations, queried directly
from the blockchain receipt log; under no attack it is exactly zero, and
under any hidden forwarding variant it rises immediately, including the two
passive variants where every cryptographic verification passes and the
proof failure indicators of Section~\ref{sec:crypto_layer} provide no
signal at all. It is therefore the only feature in the vector providing
direct signal for passive hidden forwarding independently of the
cryptographic layer, which is also why, as shown in
Algorithm~\ref{alg:lrad_rsu}, it contributes directly to Signatures S7 and
S8 rather than solely through the classifier.

The classifier is suppressed wherever a rule based signature has already
resolved the cycle, given in Equation~\eqref{eq:lstm_gate}.
\begin{equation}
	\texttt{flag}_{LSTM}(r,t) =
	\begin{cases}
		D_{LSTM}^{(k)}(t) & \text{if } \neg(flag_{S3} \lor flag_{S4}) \\
		0 & \text{if } flag_{S3} \lor flag_{S4}
	\end{cases}
	\label{eq:lstm_gate}
\end{equation}
As given in Equation~\eqref{eq:lstm_gate}, suppression is unconditional
whenever S3 or S4 fires, since residual table occupancy from those attacks
inflates reconstruction error at roadside units that are not themselves
under attack, for as long as the occupancy persists; this is a structural
effect the classifier cannot be conditioned around, so the rule based and
blockchain layers are treated as definitive for those two variants instead.

The encoder's hidden state evolves as given in
Equation~\eqref{eq:lstm_hidden}.
\begin{equation}
	\mathbf{h}_t = \textsc{LSTM}\!\left(\mathbf{x}_t, \mathbf{h}_{t-1};\, \mathbf{W}_{local}^{(k)}\right)
	\label{eq:lstm_hidden}
\end{equation}
As given in Equation~\eqref{eq:lstm_hidden}, the hidden state $\mathbf{h}_t$
updates from the current eleven feature input $\mathbf{x}_t$ and the
previous hidden state under weights $\mathbf{W}_{local}^{(k)}$ trained
locally at roadside unit $k$, so the representation carries forward
exactly the sequence structure a single window snapshot would discard.
The encoder is first trained as an autoencoder, whose reconstruction error
is the anomaly score given in Equation~\eqref{eq:anomaly_score}.
\begin{equation}
	\mathcal{A}_t^{(k)} = \frac{1}{WF}\sum_{\tau=t-W+1}^{t} \left\|\mathbf{x}_\tau - \hat{\mathbf{x}}_\tau\right\|_2^2
	\label{eq:anomaly_score}
\end{equation}
As given in Equation~\eqref{eq:anomaly_score}, this is the mean squared
reconstruction error over the ten timestep, eleven feature window. On this
reconstruction path, a benign calibrated threshold decision is given in
Equations~\eqref{eq:lstm_detection} and~\eqref{eq:lstm_threshold}.
\begin{equation}
	D_{LSTM}^{(k)}(t) =
	\begin{cases}
		1 & \mathcal{A}_t^{(k)} > \theta^{(k)} \\
		0 & \text{otherwise}
	\end{cases}
	\label{eq:lstm_detection}
\end{equation}
\begin{equation}
	\theta^{(k)} = \mu_{\mathcal{A}}^{(k)} + z_\alpha \cdot \sigma_{\mathcal{A}}^{(k)}
	\label{eq:lstm_threshold}
\end{equation}
As given in Equation~\eqref{eq:lstm_threshold}, $z_\alpha{=}3.5$ is fitted to
clear a one percent false positive rate on the held out validation split
given the limited per roadside unit calibration sample available. A warm up
adaptation guards against seed to seed mobility variance, given in
Equation~\eqref{eq:theta_adapt}.
\begin{equation}
	\theta^{(k)}_{adapted} = \max\!\left(\theta^{(k)},\; \mu_{warmup}^{(k)} + z \cdot \sigma_{warmup}^{(k)}\right)
	\label{eq:theta_adapt}
\end{equation}
As given in Equation~\eqref{eq:theta_adapt}, the maximum ensures warm up
only ever raises the threshold, never lowers it, and this adaptation applies
to the reconstruction path described above only.

The detector actually used for the results reported in
Section~\ref{sec:results} does not threshold the reconstruction error
directly. The encoder above is first trained as described, its weights are
then frozen, and a lightweight classification head is trained separately on
the frozen encoder's final hidden state, given in
Equation~\eqref{eq:lstm_cls}.
\begin{equation}
	P_{atk}^{(k)}(t) = \sigma\!\left(\mathbf{w}_2^{\top}\, \mathrm{ReLU}\!\left(\mathbf{W}_1 \mathbf{h}^{(2)}_T + \mathbf{b}_1\right) + b_2\right)
	\label{eq:lstm_cls}
\end{equation}
As given in Equation~\eqref{eq:lstm_cls}, the head is trained with binary
cross entropy against injection side ground truth, never against any
quantity computed from the input vector itself, which keeps the label
independent of the features the model consumes. Detection on this path
follows a per roadside unit threshold calibrated from that unit's own
benign traffic under live operating conditions, given in
Equation~\eqref{eq:lstm_cls_threshold}.
\begin{equation}
	D_{LSTM}^{(k)}(t) = \mathbb{1}\!\left[P_{atk}^{(k)}(t) > \theta_{cls}^{(k)}\right], \qquad \theta_{cls}^{(k)} = \max\!\left(Q_{1-\alpha}\!\left(\mathcal{P}^{(k)}_{benign}\right),\; 0.05\right)
	\label{eq:lstm_cls_threshold}
\end{equation}
As given in Equation~\eqref{eq:lstm_cls_threshold}, the threshold is the
empirical benign quantile of the roadside unit's own classifier output under
live conditions rather than a value calibrated offline, since offline
calibrated thresholds were found to transfer poorly across the
heterogeneous traffic each roadside unit observes once deployed; the warm up
adaptation of Equation~\eqref{eq:theta_adapt} is not applied on this path.
The reconstruction equations above remain in the paper because they are the
encoder's training objective and back the high confidence tier reported in
Section~\ref{sec:results}, not because they are the deployed decision rule.

Local models are aggregated under the assumption that any roadside unit may
be compromised. Submissions are first checked against their committed hash,
given in Equation~\eqref{eq:bc_model_verify}.
\begin{equation}
	\mathbb{1}_{BC}^{(k)} = \textsc{SC.VerifyModelHash}\!\left(H(\mathbf{W}_{local}^{(k)}),\; \textsc{CommittedHash}^{(k)}\right)
	\label{eq:bc_model_verify}
\end{equation}
As given in Equation~\eqref{eq:bc_model_verify}, the indicator
$\mathbb{1}_{BC}^{(k)}$ is one only when the weights a roadside unit
actually submits hash to the value it pre committed on chain, so a
roadside unit cannot submit a different update than the one it declared
in advance. Surviving submissions are then averaged, weighted and filtered as given in
Equation~\eqref{eq:fed_robust}.
\begin{equation}
	\mathbf{W}_{global}^{(t+1)} = \frac{\sum_{k=1}^{K} n_k \cdot \mathbf{W}_{local}^{(k)}(t) \cdot \mathbb{1}\!\left[d\!\left(\mathbf{W}_{local}^{(k)}, \tilde{\mathbf{W}}\right) < \gamma\right]}{\sum_{k=1}^{K} n_k \cdot \mathbb{1}\!\left[d\!\left(\mathbf{W}_{local}^{(k)}, \tilde{\mathbf{W}}\right) < \gamma\right]}
	\label{eq:fed_robust}
\end{equation}
As given in Equation~\eqref{eq:fed_robust}, $\tilde{\mathbf{W}}$ is the
coordinate wise median of submitted weights and $\gamma$ is an outlier
rejection threshold, so a roadside unit whose weights deviate beyond
$\gamma$ from the median, indicating possible gradient poisoning, is
excluded from the weighted average entirely rather than merely down
weighted, preventing a single Byzantine roadside unit from shifting the
global model arbitrarily.

Algorithm~\ref{alg:brfa_v2} summarises the full aggregation round.
\begin{algorithm}[!t]
	\caption{Byzantine Robust Federated Aggregation (BRFA-v2)}
	\label{alg:brfa_v2}
	\begin{algorithmic}[1]
		\small
		\Procedure{BRFA-v2}{$\{\mathbf{W}_{local}^{(k)}\}_{k=1}^{K}, \{n_k\}, \gamma$}
		\State $\mathcal{K}_e \gets \{k : T_{r_k} \geq T_{min}\}$ \Comment{Trust gate}
		\If{$|\mathcal{K}_e| < 2f{+}1$}
		\State \textbf{abort} \Comment{Insufficient honest participants}
		\EndIf
		\For{$k \in \mathcal{K}_e$}
		\State $\mathbb{1}_{BC}^{(k)} \gets \textsc{SC.VerifyModelHash}(\mathbf{W}^{(k)})$ \Comment{Eq.~\eqref{eq:bc_model_verify}}
		\EndFor
		\State $\tilde{\mathbf{W}} \gets \texttt{COORD\_MEDIAN}(\{\mathbf{W}^{(k)}\}_{k \in \mathcal{K}_e})$
		\For{$k \in \mathcal{K}_e$}
		\State $\omega_k \gets n_k \cdot \mathbb{1}_{BC}^{(k)} \cdot \mathbb{1}[d(\mathbf{W}^{(k)}, \tilde{\mathbf{W}}) < \gamma]$
		\EndFor
		\State $\mathbf{W}_{global} \gets \sum_k \omega_k \mathbf{W}^{(k)} / \sum_k \omega_k$ \Comment{Eq.~\eqref{eq:fed_robust}}
		\State \texttt{BC.COMMIT}$(H(\mathbf{W}_{global}))$
		\State \textbf{return} $\mathbf{W}_{global}$
		\EndProcedure
	\end{algorithmic}
\end{algorithm}
As shown in Algorithm~\ref{alg:brfa_v2}, a roadside unit must clear the trust
gate before its hash is even checked, so a roadside unit already quarantined
by Section~\ref{sec:blockchain_arch} contributes nothing to the global
model regardless of what it submits, and the round aborts outright if fewer
than $2f{+}1$ participants remain trusted, rather than silently aggregating
over a minority that could no longer be considered representative.


\subsection{Witness Based Passive Attack Detection}
\label{sec:witness}

Passive hidden forwarding survives every mechanism developed above, since
the duplicated packet is a genuine, unmodified retransmission carrying a
valid signature; detection instead depends on nearby vehicles independently
observing the same transmission. A witness raises a duplication alert
whenever it observes the same packet hash addressed to two distinct
destinations, given in Equation~\eqref{eq:dup_alert_cond}.
\begin{equation}
	\text{DA}_w(p){=}1\iff \exists\,dst{\neq}dst': H(p)\in\mathcal{L}_w(dst,W) \land H(p)\in\mathcal{L}_w(dst',W)
	\label{eq:dup_alert_cond}
\end{equation}
As given in Equation~\eqref{eq:dup_alert_cond}, $\mathcal{L}_w(dst,W)$ is the
set of hashes the witness observed directed to $dst$ within window $W$; this
condition is impossible for legitimate, single path traffic, so any
occurrence is itself evidence of duplication regardless of which node is
ultimately found responsible. A related, severe case is covered separately:
a forwarding delay long enough to exceed the safety deadline $T_{fwd}$
renders a packet functionally equivalent to non delivery, given in
Equation~\eqref{eq:nfwd_detect}.
\begin{equation}
	\text{NFA}_w(p){=}1\iff H(p)\in\mathcal{L}_w^{recv}(v_i,W) \land H(p)\notin\mathcal{L}_w^{fwd}(v_i,T_{fwd})
	\label{eq:nfwd_detect}
\end{equation}
As given in Equation~\eqref{eq:nfwd_detect}, this condition holds when a
packet a witness saw arrive at $v_i$ was never witnessed leaving $v_i$
within the safety deadline, which this paper treats as a severe case of the
selective time delay variants PHANTOM addresses rather than as a ninth
attack type in its own right \cite{Phantom2026}. Both alert types are signed
before submission to the blockchain, given in Equations~\eqref{eq:da_sign}
and~\eqref{eq:nfa_sign}.
\begin{equation}
	\alpha_w = \textsc{ML-DSA-87.Sign}\!\left(sk_w,\; H_{\text{SHA3-512}}\!\left(H(p) \,\|\, dst \,\|\, dst' \,\|\, ts_w\right)\right)
	\label{eq:da_sign}
\end{equation}
\begin{equation}
	\beta_w = \textsc{ML-DSA-87.Sign}\!\left(sk_w,\; H_{\text{SHA3-512}}\!\left(H(p) \,\|\, v_i \,\|\, ts_w \,\|\, T_{fwd}\right)\right)
	\label{eq:nfa_sign}
\end{equation}
As given in Equations~\eqref{eq:da_sign} and~\eqref{eq:nfa_sign}, lattice
based signing rather than a symmetric tag is required here specifically
because both alerts are written directly to the blockchain and must be
attributable to a specific witness vehicle by any third party verifier
without that verifier holding a pre shared key.

A single witness alert must never be sufficient to penalize a node, since
one compromised or mistaken witness could otherwise accuse an innocent node
at will. The smart contract penalty trigger requires a distinct witness
quorum instead, given in Equation~\eqref{eq:bft_penalty}.
\begin{equation}
	\textsc{SC.PenalizeTrust}(v_i) \Leftarrow \Bigl|\bigl\{ w : \exists\,(\alpha_w \text{ or } \beta_w) \text{ s.t. } \textsc{ML-DSA-87.Verify}(\cdot,pk_w,\cdot) = 1 \;\land\; \text{event}(v_i,p,w)\text{ confirmed} \bigr\}\Bigr| \geq 2f{+}1
	\label{eq:bft_penalty}
\end{equation}
As given in Equation~\eqref{eq:bft_penalty}, the cardinality is taken over
distinct witness vehicles that each independently submitted a verified
alert for the same event, not over the number of alert messages, so a
single vehicle resubmitting its own alert can never cross the threshold
unilaterally. Once a node is quarantined by this mechanism or by
Section~\ref{sec:blockchain_arch}, Signatures S5 through S8 stop asserting
against it, since a detector that keeps accusing a node the system has
already contained measures its own failure to switch off rather than a
genuine gap in detection, an exception encoded directly in
Algorithm~\ref{alg:lrad_rsu} above.


\subsection{Blockchain Architecture}
\label{sec:blockchain_arch}

All blockchain writes originate exclusively from roadside units, and for
this paper's own signatures the packet reception log is the single most
load bearing mechanism, since it is what makes cross roadside unit
duplication checking possible without a direct communication channel
between the two roadside units involved. Figure~\ref{fig:blockchain_arch}
shows the resulting two tier structure.
\begin{figure}[!t]
	\centering
	\includegraphics[width=0.95\linewidth]{Figure_4_Final.eps}
	\caption{Blockchain Architecture and Trust Management}
	\label{fig:blockchain_arch}
\end{figure}

A flow modification is installed only once endorsed by $f{+}1$ independent
roadside units, given in Equations~\eqref{eq:rsu_endorsement}
and~\eqref{eq:endorsed_commit}.
\begin{equation}
	\epsilon_j = \textsc{ML-DSA-87.Sign}\!\left(sk_{r_j},\; H_{\text{SHA3-512}}\!\left(\texttt{FlowMod} \,\|\, ts \,\|\, \mathcal{T}_{r_j}\right)\right)
	\label{eq:rsu_endorsement}
\end{equation}
\begin{equation}
	C_{\mathcal{P}} = \textsc{BC.Commit}\!\left(H(\texttt{FlowMod}) \,\|\, \{\epsilon_j\}_{j=1}^{f+1} \,\|\, ts_{commit}\right)
	\label{eq:endorsed_commit}
\end{equation}
As given in Equations~\eqref{eq:rsu_endorsement}
and~\eqref{eq:endorsed_commit}, a controller issuing an unauthorized flow
modification toward an unauthorized destination cannot collect $f{+}1$
honest endorsements and so cannot produce a valid commitment, which is
exactly the absence Signature S5's second conjunct tests for. Every flow
modification received is logged before installation, given in
Equation~\eqref{eq:flowmod_log}, and every packet reception is logged
separately, given in Equation~\eqref{eq:packet_receipt_log}.
\begin{equation}
	\textsc{BC.LogFlowMod}\!\left(r_k,\; H(\texttt{FlowMod}_{recv}) \,\|\, ts_{recv} \,\|\, \epsilon_{endorsement}\right)
	\label{eq:flowmod_log}
\end{equation}
\begin{equation}
	\textsc{BC.LogReceipt}\!\left(H(p),\; dst,\; r_k,\; ts_{recv},\; \textsc{ML-DSA-87.Verify}(\sigma_{copy},pk_s,m_{copy})\right)
	\label{eq:packet_receipt_log}
\end{equation}
As given in Equation~\eqref{eq:packet_receipt_log}, any roadside unit may
query whether a given packet hash was received at a specific destination
within a window through $\textsc{BC.QueryReceipt}(H(p), d', W)$, which is
the query every one of this paper's signatures and the reception log
anomaly feature ultimately rest on. Absence of a corresponding commitment
flags a flow modification as unauthorized, given in
Equation~\eqref{eq:unauth_flowmod}.
\begin{equation}
	f_{unauth}(r_k, t) =
	\begin{cases}
		1 & \texttt{FlowMod}_{recv} \notin \textsc{BC.Query}(\mathcal{C}_{\mathcal{P}}) \\
		0 & \text{otherwise}
	\end{cases}
	\label{eq:unauth_flowmod}
\end{equation}
As given in Equation~\eqref{eq:unauth_flowmod}, this flag is a
pre installation enforcement gate rather than only a detection signal: a
roadside unit evaluating it as one rejects the flow modification outright,
so the underlying rule never takes effect regardless of what any signature
later concludes about it.

Node level trust updates on every cryptographic verification outcome, given
in Equation~\eqref{eq:trust_update}, and quarantine follows automatically
once trust drops below the minimum, given in Equation~\eqref{eq:quarantine}.
\begin{equation}
	T_v^{(t+1)} =
	\begin{cases}
		\min(1,\, T_v^{(t)} + \Delta_r) & \textsc{BatchVerify} = 1 \;\land\; b_{\pi} = 1 \\
		\max(0,\, T_v^{(t)} - \Delta_p) & \textsc{BatchVerify} = 0 \;\lor\; b_{\pi} = 0
	\end{cases}
	\label{eq:trust_update}
\end{equation}
\begin{equation}
	\textsc{SC.Quarantine}(v) \Leftarrow T_v^{(t)} < T_{min}
	\label{eq:quarantine}
\end{equation}
As given in Equation~\eqref{eq:quarantine}, quarantine triggers
automatically once trust crosses the threshold, without a centralized
decision. Enforcement of quarantine, meaning whether a quarantined node is
actually blocked from continuing to schedule a hidden duplicate rather than
merely being flagged, is an explicit configuration switch in this project:
the ablation results in Section~\ref{sec:results} are measured with
enforcement disabled so that they isolate detection quality on its own
terms, and its effect once enabled is evaluated separately, a distinction
this paper states plainly since it materially changes what a quarantine
count in any table is allowed to claim. Finally, the authorized forwarding
policy $\mathcal{P}(s,d)$ that Signature S5's first conjunct tests against
is itself committed by the endorsing roadside units at flow setup, given in
Equation~\eqref{eq:policy_commit}.
\begin{equation}
	C_{\mathcal{P}} = \textsc{Commit}\!\left(H(\mathcal{P}(s,d)) \,\|\, ts_{setup} \,\|\, \{\epsilon_j\}_{j=1}^{f+1}\right)
	\label{eq:policy_commit}
\end{equation}
As given in Equation~\eqref{eq:policy_commit}, this commitment is what
Signature S5's first conjunct ultimately verifies a proposed duplicate
destination against, so the authorized policy a flow is held to is never
dependent on the controller's own, potentially compromised, account of
what it authorized, closing the same gap at the policy layer that the
addressing conjunct closes at the signature layer.



\section{Evaluation Setup}
\label{sec:evaluation_setup}

All five experiments in Section~\ref{sec:results} evaluate Signatures S5
through S8 exclusively, against the same three baselines named in
Section~\ref{sec:related}, of which only FADE addresses this paper's own
attack class. Table~\ref{tab:vanguard_sim_settings} lists the settings
required for replication; parameters that are not themselves part of this
paper's design, together with their full sensitivity sweeps, are reported
in the companion papers. Trip generation across the shared mobility trace
follows a high random routing factor to avoid unnatural vehicle
clustering \cite{sumo_randomtrips}. The DSRC transmission power is set
within the standard power class the Federal Communications Commission
allocated for this band \cite{FCC03-324}, and the routing update and data
transmission clock is set below the Basic Safety Message rate the SAE
J2735 standard specifies \cite{sae_j2735}, so as to remain frequent
enough to observe the signatures of Section~\ref{system_model} within
their own observation window. The reconstruction path threshold's
$z_\alpha{=}3.5$ setting of Equation~\eqref{eq:lstm_threshold} follows the
same three sigma convention for bounding false positive rate under a
Gaussian assumption \cite{pukelsheim1994three}, distinct from the
empirical quantile the deployed classifier threshold actually uses.

\begin{table}[!t]
	\centering
	\caption{Simulation Settings}
	\label{tab:vanguard_sim_settings}
	\renewcommand{\arraystretch}{1.2}
	\footnotesize
	\begin{tabular}{|>{\raggedright\arraybackslash}p{3.2cm}|>{\raggedright\arraybackslash}p{4.4cm}|}
		\hline
		\textbf{Parameter} & \textbf{Value} \\
		\hline
		Mobility trace & Los Angeles, California (OpenStreetMap, SUMO) \\
		\hline
		Network scale (default) & 200 vehicles, 64 RSUs (8$\times$8 grid), 4 controllers \\
		\hline
		Data / control plane & IEEE 802.11p DSRC / Ethernet CSMA \\
		\hline
		LSTM classifier threshold & Per RSU empirical benign quantile, $\alpha = 0.01$, floor $0.05$ \\
		\hline
		LSTM input, HF features & Eleven feature vector; destination diversity, throughput asymmetry, and reception anomaly enter as $\log(1{+}x)$ \\
		\hline
		Byzantine aggregation & Krum threshold $\gamma \in \{1.5, 2.0, 2.5, 3.0\}$ times the coordinate wise standard deviation of submitted weights, swept for best poisoning resistance under a 25\% malicious RSU fraction; trust gate $T_{min} \in \{0.3, 0.5, 0.7\}$, shared with the blockchain quarantine threshold of Equation~\eqref{eq:quarantine}; neither is finalized in this draft \\
		\hline
		Cryptographic layer & ML-DSA-87 signatures; batch size $B \in \{10, 15, 20\}$ packets/batch, swept for lowest end to end latency within the 50\,ms bound, not yet finalized in this draft \\
		\hline
		Blockchain endorsement & $f{+}1$ RSU endorsements per FlowMod; $f$ is systematically varied against simulated adversarial node density as part of Byzantine tolerance testing rather than held at one fixed value \\
		\hline
		Witness quorum & $2f{+}1$ distinct witnesses, with $f{=}1$ (quorum of 3) expected under mean vehicle density, though the same variation applies here as to the endorsement parameter above; window $W{=}10$s \\
		\hline
		Duplication / volume windows & S6 duplication window 30s; S7 tumbling volume window 10s \\
		\hline
		Baselines & TAP, SFTO Guard, FADE \\
		\hline
		Ground truth & Injection side, independent of any detector input feature \\
		\hline
	\end{tabular}
\end{table}

Detection quality is reported as MCC for Signatures S5 through S8
individually and pooled, following the same evaluation methodology as the
companion papers: per RSU, ten second deduplicated blocks, suspect rather
than observer attribution, run length always stated, and hidden forwarding
ground truth latched from a node's first attack action until quarantine
rather than for the run's duration, since the persistent, ongoing nature of
a compromised routing rule is what Signatures S5 through S8 actually test
for. Unauthorized copy rate, M3, is the share of hidden forwarding
duplicates that reach an unauthorized destination undetected. Witness alert
precision and recall, M12, is reported separately from the composite window
score wherever the witness mechanism's own internal counters are the more
direct measurement, since a window level score driven by
Equation~\eqref{eq:bft_penalty}'s quorum requirement can otherwise
understate a witness path that is working correctly but has not yet reached
quorum within a given block. Mitigation latency, M4, and end to end
latency, M6, are defined identically to the companion papers.


\section{Results}
\label{sec:results}

% All quantitative values in this section are placeholders pending the
% simulation runs specified in each experiment below, executed under the
% finalised, debugged configuration of Table~\ref{tab:vanguard_sim_settings}.
% The analysis prose states the mechanism each trend is expected to follow
% and should be checked against, not discarded, once real numbers land.
% Experiment 4 specifically must not be populated until it is confirmed
% whether its underlying data collection predates or postdates the
% destination diversity feature correction described in
% Section~\ref{sec:fed_lstm}; this is a standing, unresolved blocker, not a
% formatting note.

\subsection{Experiment 1: Detection Under Varying Attack Penetration and Forwarding Intensity}

Attack penetration is swept over $p \in \{0,20,40,60,80,100\}\%$ at three
forwarding intensity levels, $\{25\%, 50\%, 100\%\}$ of eligible flows, for
Signatures S5 through S8 only.
\begin{figure}[!t]
	\centering
	\fbox{\begin{minipage}[c][4.2cm]{0.92\linewidth}
			\centering
			Figure placeholder: macro MCC versus attack penetration, one line
			per forwarding intensity, VANGUARD HF against TAP, SFTO Guard,
			and FADE.
	\end{minipage}}
	\caption{Detection Quality Under Varying Penetration and Forwarding Intensity}
	\label{fig:exp1_mcc}
\end{figure}
Active variants, S5 and S6, are expected to show comparatively little
sensitivity to forwarding intensity, since aggregate signature verification
either fails or it does not: a single fabricated duplicate is as
cryptographically detectable as many, so detection quality for these two
variants should be governed more by penetration than by intensity. Passive
variants, S7 and S8, are expected to show the opposite pattern, since the
reception log anomaly score of Equation~\eqref{eq:feat_ranom} accumulates
roughly in proportion to how many unauthorized receptions occur per unit
time, so higher forwarding intensity should improve passive detection
directly rather than only through more attack instances to detect. FADE is
expected to compete at high intensity, where duplication is frequent enough
to disturb its flow conservation test, and to degrade sharply at low
intensity, where the same test has too few samples within a vehicle's zone
residence time to be reliable, exactly the limitation
Section~\ref{sec:dual_mode} identifies. At $p{=}\TBD\%$ and intensity
$\TBD\%$, macro MCC is $\TBD$ against $\TBD$ for the strongest baseline.

\subsection{Experiment 2: Detection Across Vehicular Speed}

Mean vehicle speed is swept over $\{10, 60, 100, 140\}$km/h, reusing the
same SUMO mobility traces generated for the companion paper PHANTOM where
the logged features cover this paper's requirements as well, avoiding a
redundant mobility simulation for the same environment \cite{Phantom2026}.
\begin{figure}[!t]
	\centering
	\fbox{\begin{minipage}[c][4.2cm]{0.92\linewidth}
			\centering
			Figure placeholder: macro MCC versus mean vehicle speed, split by
			active (S5, S6) and passive (S7, S8) variants, against the three
			baselines.
	\end{minipage}}
	\caption{Detection Quality Across Vehicular Speed}
	\label{fig:exp2_speed}
\end{figure}
Unlike PHANTOM's own speed experiment, this paper's primary detectors do not
depend on a mobility adjusted timing threshold at all, so active variant
detection is expected to be largely speed independent, governed by
cryptographic verification rather than by any quantity that scales with
speed. Passive variant detection is expected to be more sensitive to speed
than active detection, but for a different reason than timing: the witness
quorum of Equation~\eqref{eq:bft_penalty} requires $2f{+}1$ distinct
vehicles to remain within range long enough to observe and confirm the same
event, and higher speed shortens how long any given set of neighbouring
vehicles remains co located, which could plausibly slow how quickly quorum
is reached even if it does not prevent it outright. At
$\bar{v}{=}140$km/h, macro MCC for the active variants is $\TBD$ and for
the passive variants is $\TBD$.

\subsection{Experiment 3: Scalability}

Network scale is varied as $N_v \in \{100,200,300,400\}$ vehicles, again
reusing PHANTOM's mobility runs where possible, holding the roadside unit
grid and controller count fixed, for Signatures S5 through S8 only.
\begin{figure}[!t]
	\centering
	\fbox{\begin{minipage}[c][4.2cm]{0.92\linewidth}
			\centering
			Figure placeholder: macro MCC and end to end latency versus
			vehicle count, VANGUARD HF against the three baselines.
	\end{minipage}}
	\caption{Scalability of Detection Quality and Latency}
	\label{fig:exp3_scale}
\end{figure}
Detection quality is expected to hold steady as vehicle count grows, since
per roadside unit classifier calibration of
Equation~\eqref{eq:lstm_cls_threshold} depends only on that unit's own
benign traffic, not on total network size. The Byzantine robust aggregation
round of Algorithm~\ref{alg:brfa_v2} scales with roadside unit count, which
is held fixed here, so its cost should not grow with vehicle count either,
though this is the binding cost as the network grows along the roadside
unit axis instead, a distinction developed further in
Section~\ref{sec:discussion}. At $N_v{=}400$, macro MCC is $\TBD$ and end to
end latency is $\TBD$ms against the hundred millisecond bound.

\subsection{Experiment 4: Detection Under Decreasing Attack Observable Evidence}

\textit{This experiment's results are withheld pending confirmation of
whether its underlying data collection predates or postdates the
destination diversity feature correction of
Section~\ref{sec:fed_lstm}. The experimental design below is final; the
numbers are not, and should not be filled in from any existing run without
that confirmation first.}

Attack observable evidence is swept over four points combining a
selectivity ratio with copy destination divergence: $\text{AOEI}{=}0.25$ at
10\% targeting and $d_{div}{=}0.8 d_{max}$, $\text{AOEI}{=}0.50$ at 25\%
targeting and $d_{div}{=}0.5 d_{max}$, $\text{AOEI}{=}0.75$ at 75\%
targeting and $d_{div}{=}0.2 d_{max}$, and $\text{AOEI}{=}1.0$ at 100\%
targeting and $d_{div}{=}0$.
\begin{figure}[!t]
	\centering
	\fbox{\begin{minipage}[c][4.2cm]{0.92\linewidth}
			\centering
			Figure placeholder: macro MCC and unauthorized copy rate versus
			attack observable evidence index, VANGUARD HF against FADE.
	\end{minipage}}
	\caption{Detection Quality Under Decreasing Observable Evidence}
	\label{fig:exp4_aoei}
\end{figure}
As copy destination divergence falls toward zero, the unauthorized
destination moves closer to the authorized one; this does not change
whether Equation~\eqref{eq:feat_ddiv}'s policy membership test flags the
destination as unauthorized, since that test is categorical rather than
distance based, but it may plausibly affect how far a witness vehicle near
the authorized destination has to be for the duplicate to fall within its
own observation range, which would connect this experiment to witness
quorum feasibility rather than to the classifier at all. This mechanism is
stated as a hypothesis to check, not a finding, given the open data
provenance question above. At $\text{AOEI}{=}0.25$, macro MCC is $\TBD$
against $\TBD$ at $\text{AOEI}{=}1.0$.

\subsection{Experiment 5: Per Variant Comparison Under Default Settings}

Table~\ref{tab:exp5_variant} reports macro MCC, unauthorized copy rate,
mitigation latency, and end to end latency for each of Signatures S5
through S8 individually, under the default settings of
Table~\ref{tab:vanguard_sim_settings}. TAP and SFTO Guard address selective
time delay and slow flow table exhaustion respectively and are not
applicable to any variant reported here.
\begin{table}[!t]
	\centering
	\caption{Per Variant Detection Quality, Default Settings}
	\label{tab:exp5_variant}
	\renewcommand{\arraystretch}{1.15}
	\footnotesize
	\begin{tabular}{|c|c|c|c|c|}
		\hline
		\textbf{Variant} & \textbf{VANGUARD HF} & \textbf{FADE} & \textbf{M3} & \textbf{M4} \\
		\hline
		S5 & \TBD & \TBD & \TBD\% & \TBD \\
		\hline
		S6 & \TBD & \TBD & \TBD\% & \TBD \\
		\hline
		S7 & \TBD & \TBD & \TBD\% & \TBD \\
		\hline
		S8 & \TBD & \TBD & \TBD\% & \TBD \\
		\hline
	\end{tabular}
\end{table}
FADE is expected to compete most closely on the active variants, S5 and S6,
where a fabricated duplicate still disturbs the flow conservation signal
FADE was designed around, and to fall furthest behind on the passive
variants, S7 and S8, where the retransmitted packet is genuinely authentic
and disturbs nothing FADE observes at all. This is the central expected
result of this experiment: a baseline built for one hidden forwarding
mechanism should not transfer to the other, which is exactly why this
paper's own witness layer, rather than any single cryptographic or
statistical check, is necessary to cover both.


\section{Discussion}
\label{sec:discussion}

Mitigation latency is reported throughout this paper as a pair rather than
a single mean: a prevention rate, the share of quarantined attackers
stopped before their first attack action, and a mean latency computed only
over attackers that acted before containment, since a single mean over both
populations moves in the wrong direction as containment improves.

\subsection{Statistical Significance}
Every experiment in Section~\ref{sec:results} is run across five seeds. The
comparison between VANGUARD HF and FADE, the one baseline that addresses
this paper's own attack class at all, is tested with a paired test across
seeds in Experiment 1 rather than a comparison of seed averages, since both
methods share the same mobility realisations. At $p{=}\TBD\%$ and intensity
$\TBD\%$, the macro MCC difference is $\TBD$ with $\TBD$ ($n{=}5$ paired
seeds); this is reported once the experiments are complete.

\subsection{Strengths and Weaknesses}
The central strength this paper demonstrates is coverage of passive hidden
forwarding, a case no baseline in Section~\ref{sec:related} addresses at
all, since a passively retransmitted packet is a genuinely authentic copy
that satisfies every cryptographic check available. The addressing
conjunct added to Signatures S5 through S8 is what makes these signatures
well posed on a broadcast medium in the first place, without it every
neighbour that merely overhears a transmission satisfies the remaining
conjuncts equally, and the witness quorum of
Equation~\eqref{eq:bft_penalty} then covers exactly the residual case those
signatures cannot: detection with no cryptographic or volumetric trace to
observe at all.

Three weaknesses are worth stating plainly. First, as in the companion
papers, the composite decision across Signatures S5 through S8 and the
classifier is a logical OR, trading false positive rate for coverage.
Second, passive detection depends on witness density: the $2f{+}1$ quorum
becomes harder to satisfy exactly where fewer vehicles are present to
observe an event, a dependency this paper's architecture does not remove,
only bounds through the Byzantine fault tolerance parameter. Third,
Experiment 4 of Section~\ref{sec:results} remains withheld pending
confirmation of its data provenance relative to the destination diversity
feature correction of Section~\ref{sec:fed_lstm}; this paper does not
report a number for that experiment rather than risk reporting one built on
the wrong feature definition.

\subsection{Runtime and Space Complexity}
The Byzantine robust aggregation round of Algorithm~\ref{alg:brfa_v2} is
quadratic in roadside unit count, and per hop signature verification is
linear in path length. Appendix~\ref{appendix} states and proves the formal
bounds for both, since this paper carries the full classifier and
aggregation architecture; the timing compliance proof cost specific to
selective time delay is developed in the companion paper PHANTOM
\cite{Phantom2026}.

\subsection{Scalability}
Detection quality is expected to hold steady as vehicle count grows, since
per roadside unit classifier calibration depends only on that unit's own
benign traffic, consistent with Experiment 3 of Section~\ref{sec:results}.
The aggregation round above is the binding cost as the network grows along
the roadside unit axis specifically, not the vehicle axis Experiment 3
varies.

\subsection{Deployment Feasibility}
Federated training keeps raw traffic local to each roadside unit, sharing
only model weights, which is a genuine privacy property rather than an
incidental one given this paper's attack class is itself about unauthorized
observation of traffic. The witness mechanism requires no additional
hardware beyond a vehicle's existing radio, since it depends only on
passively overhearing a broadcast transmission already occurring on the
shared medium.

\subsection{Convergence}
Byzantine robust aggregation converges toward the honest aggregate under
the bounded malicious fraction assumed throughout this project.
Appendix~\ref{appendix} states and proves this bound in full, since this
paper carries the complete classifier and aggregation architecture the
bound applies to; the companion paper PHANTOM refers to this appendix
rather than repeating the proof for its own, abbreviated treatment of the
same classifier \cite{Phantom2026}.

\subsection{Security Analysis}
Falsely accusing a benign node under Signatures S5 through S8 requires
either breaking ML-DSA-87 unforgeability or compromising the physical link
layer addressing the addressed conjunct is read from, and the witness
quorum of Equation~\eqref{eq:bft_penalty} adds a second, independent
barrier requiring collusion across $2f{+}1$ witnesses on top of that.
Appendix~\ref{appendix} states and proves both claims in full. The
controller trust channel's own Byzantine fault tolerance, keyed to
Signature S1 in the companion paper, is proved there rather than repeated
here, since it is not keyed to any signature this paper introduces
\cite{Phantom2026}.


\section{Conclusion and Future Work}

Cryptographic verification and flow conservation testing each fail against
part of the hidden forwarding threat: the former cannot reject a passively
retransmitted, genuinely authentic packet, and the latter needs more
samples than a vehicle's zone residence time reliably provides. VANGUARD HF
addresses both failures with signatures corrected to be well posed on a
broadcast medium, a federated classifier that provides direct signal for
the passive case independently of any cryptographic observable, and a
witness quorum mechanism that covers what neither the classifier nor
cryptography alone can. Across the five experiments of
Section~\ref{sec:results}, VANGUARD HF is shown to sustain macro MCC of
$\TBD$ against $\TBD$ for FADE, the one baseline addressing this attack
class, at default settings. The main limitation is Experiment 4, withheld
pending confirmation of a data provenance question this paper states
plainly rather than resolves by assumption.

Two companion papers extend this evaluation: PHANTOM develops the
selective time delay signatures and the mobility adjusted threshold this
paper does not address \cite{Phantom2026}, and NEXUS reports both attack
classes evaluated jointly under simultaneous attack \cite{Nexus2026}.
Future work includes resolving Experiment 4's data provenance question and
re running it once settled, a density adaptive witness quorum that lowers
the Byzantine fault tolerance parameter where vehicle density is too sparse
to sustain $2f{+}1$ distinct witnesses, and replacing the reception log
anomaly score's binary contribution to Signatures S7 and S8 with a graded
one that reflects its magnitude rather than only whether it is nonzero.

\subsection{Acknowledgments}
\noindent None

\subsection{Conflict of interest}
\noindent Authors declare no conflict of interest.

\section[\appendixname~\thesection]{Appendix}
\label{appendix}

This appendix gives the complexity, convergence, and security analysis for
this paper's own contributions.


\subsection{Complexity Analysis}
\label{app:complexity}

Table~\ref{tab:vanguard_complexity} states the per operation cost of the
full classifier and aggregation architecture this paper carries, on top of
the aggregate signature verification shared with the companion papers.
\begin{table}[!t]
	\centering
	\caption{Runtime and Space Complexity, Classifier and Aggregation}
	\label{tab:vanguard_complexity}
	\renewcommand{\arraystretch}{1.2}
	\footnotesize
	\begin{tabular}{|l|c|c|}
		\hline
		\textbf{Operation} & \textbf{Time} & \textbf{Space} \\
		\hline
		BatchVerify, Eq.~\eqref{eq:batch_verify} & $O(n)$ & $O(n)$ \\
		\hline
		STARK.Prove/Verify, $\pi_{hop}$ & $O(1)$ per hop & $O(1)$ per hop \\
		\hline
		Classifier inference, Eq.~\eqref{eq:lstm_cls} & $O(1)$ per RSU per cycle & $O(1)$ per RSU \\
		\hline
		Krum filter, Algorithm~\ref{alg:brfa_v2} & $O(K^2)$ & $O(Kd)$ \\
		\hline
		Endorsement, Section~\ref{sec:blockchain_arch} & $O(f{+}1)$ & $O(f{+}1)$ \\
		\hline
	\end{tabular}
\end{table}
As given in Table~\ref{tab:vanguard_complexity}, $n$ is path length in hops,
$K$ the number of participating roadside units, $d$ the model weight
dimension, and $f$ the Byzantine fault parameter.

\begin{lemma}
	\label{lem:vanguard_complexity}
	The per packet verification cost of the cryptographic layer is $O(n)$
	in path length $n$, and the per round cost of Byzantine robust
	aggregation is $O(K^2)$ in the number of participating roadside units
	$K$, independently of $n$.
\end{lemma}
\begin{proof}
	BatchVerify of Equation~\eqref{eq:batch_verify} accumulates $n$ per hop
	signatures into a single combined check, which is $O(n)$; the hop
	legitimacy proof $\pi_{hop}$ is generated and verified once per hop at
	cost independent of $n$ beyond the single proof already accounted for
	at that hop, so it does not change this order. Classifier inference at
	a given roadside unit and cycle is a single forward pass over a fixed
	size input and a fixed size hidden state, hence $O(1)$ per roadside
	unit per cycle regardless of $n$ or $K$. The Krum filter of
	Algorithm~\ref{alg:brfa_v2} computes a pairwise distance between every
	pair of the $K$ submitted models to identify the coordinate wise
	median neighbourhood, which requires $\binom{K}{2}$ distance
	computations and is therefore $O(K^2)$, independent of $n$.
	\qed
\end{proof}

\begin{corollary}
	At the default network scale of Table~\ref{tab:vanguard_sim_settings},
	with $K{=}64$ roadside units, the Krum filter dominates the total per
	round cost whenever the evaluated path length satisfies $n \ll K^2$,
	which holds for every path length reachable within the eight by eight
	grid deployment.
\end{corollary}

\begin{remark}[Limitation]
	The quadratic term in Lemma~\ref{lem:vanguard_complexity} is not
	exercised by Experiment 3 of Section~\ref{sec:results}, which scales
	vehicle count rather than roadside unit count; a deployment that
	instead scales the roadside unit grid would need the aggregation round
	cost re evaluated at the larger $K$. The timing compliance proof
	specific to selective time delay, and its associated cost, belongs to
	the companion paper PHANTOM and is not included in
	Table~\ref{tab:vanguard_complexity} \cite{Phantom2026}.
\end{remark}

\subsection{Convergence Analysis}
\label{app:convergence}

\begin{lemma}
	\label{lem:convergence}
	Let $\rho_{mal}$ be the fraction of roadside units submitting
	adversarial weight updates in a given aggregation round. If
	$\rho_{mal} < 1/3$, the output of Algorithm~\ref{alg:brfa_v2} lies
	within a bounded distance of the aggregate that would be produced by
	the honest submissions alone, for a Krum threshold $\gamma$ chosen at
	least as large as the spread of the honest population.
\end{lemma}
\begin{proof}
	The trust gate and hash verification steps of Algorithm~\ref{alg:brfa_v2}
	admit only submissions from roadside units above the operational trust
	threshold with a weight matching their committed hash, given in
	Equation~\eqref{eq:bc_model_verify}, so any adversarial submission
	reaching the Krum filter carries no advantage from forged identity or
	a mismatched payload. The coordinate wise median $\tilde{\mathbf{W}}$
	used by the filter is itself robust to up to, but not including, half
	of the remaining population being adversarial: each coordinate of the
	median is bounded between the corresponding coordinates of the
	honest submissions regardless of how extreme the adversarial values
	are, provided honest submissions remain the majority. For
	$\rho_{mal}<1/3$, the honest population is comfortably the majority
	required, so the median stays within the convex spread of the honest
	submissions. Any submission further than $\gamma$ from
	$\tilde{\mathbf{W}}$ is excluded before averaging in
	Equation~\eqref{eq:fed_robust}; once $\gamma$ is set at least as large
	as the honest population's own spread around the median, every honest
	submission is retained while an adversarial submission distant enough
	to bias the average by more than the stated bound is excluded, so the
	weighted aggregation step of Algorithm~\ref{alg:brfa_v2} stays within
	a bounded distance of the honest only aggregate.
	\qed
\end{proof}

\begin{remark}[Limitation]
	Lemma~\ref{lem:convergence} bounds the deviation from the honest
	aggregate per round rather than guaranteeing convergence of the
	underlying model to a global optimum, which additionally depends on
	the non convexity of the encoder trained in
	Section~\ref{sec:fed_lstm}. Empirical support for the stated bound
	under this project's Byzantine robust aggregation ablation will be
	reported alongside this paper's own results rather than asserted here
	without measurement.
\end{remark}

\subsection{Security Analysis}
\label{app:security}

\begin{lemma}
	\label{lem:s5_unforgeability}
	Falsely accusing a benign forwarding node under Signatures S5 through
	S8 requires either breaking ML-DSA-87 unforgeability or compromising
	the link layer addressing that the addressed conjunct is read from.
\end{lemma}
\begin{proof}
	A benign forwarding node relays a packet under a signature chain
	originating from its genuine sender, so
	$\textsc{ML-DSA-87.Verify}(\sigma_{copy},pk_s,m_{copy})=1$ for any copy
	the accusing node genuinely received. Producing a copy that fails this
	check while attributing it to a benign node's forwarding action
	requires forging a valid looking signature under the sender's public
	key, precluded by the module lattice hardness assumption underlying
	ML-DSA-87 unforgeability. For Signatures S7 and S8, where verification
	succeeding rather than failing is itself part of the condition, the
	same argument applies in reverse: a benign node cannot be made to
	appear as though it fabricated a duplicate, since fabrication requires
	the copy to fail verification, which a genuinely retransmitted,
	authentically signed packet does not. The addressing conjunct
	$\textsc{Addressed}(p,d')=1$, common to all four signatures, is read
	directly from the link layer destination of the frame the accusing
	node itself received, independent of the signed next hop, so
	satisfying it without genuinely having been addressed the duplicate
	requires control over the physical radio channel rather than over any
	cryptographic or protocol state.
	\qed
\end{proof}

\begin{corollary}
	Since Equation~\eqref{eq:bft_penalty} additionally requires $2f{+}1$
	distinct witnesses to independently confirm an event before a trust
	penalty is applied through the witness path, an adversary must both
	defeat Lemma~\ref{lem:s5_unforgeability} and collude across $2f{+}1$
	witnesses to falsely penalise a node through that path, exceeding the
	Byzantine bound assumed throughout this project.
\end{corollary}

\begin{remark}[Limitation]
	Lemma~\ref{lem:s5_unforgeability} does not, by itself, explain why a
	genuinely malicious passive duplication is detected at all, only why a
	benign node cannot be falsely accused: a passively retransmitted
	packet is a genuinely authentic copy that satisfies every
	cryptographic check by construction for both the accuser and, had it
	been legitimate, would satisfy them for a genuine sender too. Coverage
	of genuinely malicious passive duplication rests on the witness
	quorum of Section~\ref{sec:witness} observing the duplication
	directly, not on any cryptographic hardness assumption, which is
	precisely why the witness mechanism exists as an independent detection
	path rather than as a corroborating check on the signatures above.
\end{remark}

\begin{remark}
	The controller trust channel's own Byzantine fault tolerance is proved
	in the companion paper PHANTOM, since that channel is keyed to
	Signature S1, a signature this paper does not introduce, and is not
	repeated here \cite{Phantom2026}.
\end{remark}

\bibliographystyle{elsarticle-num}
\begin{thebibliography}{00}

	% NOTE: entries marked [verified] were checked against publicly indexed
	% records this session. Entries marked [confirm before submission] rest on
	% recollection from earlier in this project and must be checked against
	% the publisher record before submission, per the reference integrity
	% standard already applied to the companion papers.
	%
	% CORRECTION LOG (this session): the FADE and FTISCON entries below were
	% previously attributed to the wrong authors and venues (Zhang2021FADE,
	% Krishnan2023). Both papers are real; the attribution was not. Corrected
	% against the actual publisher record, not removed, since the underlying
	% work exists and is genuinely relevant.

	\bibitem{Mohammadi2016} % [verified]
	A.~A. Mohammadi, P. Kazemian, M. Pakravan, Detecting malicious packet drops and misroutings using header space analysis, in: Proc. 8th Int. Symp. Telecommun. (IST), Tehran, Iran, 2016, pp. 521--526.

	\bibitem{Pang2016FADE} % [verified; corrected from Zhang2021FADE]
	C. Pang, Y. Jiang, Q. Li, FADE: detecting forwarding anomaly in software defined networks, in: Proc. IEEE Int. Conf. Commun. (ICC), 2016.

	\bibitem{Sasaki2016SDNsec} % [verified]
	T. Sasaki, C. Pappas, T. Lee, T. Hoefler, A. Perrig, SDNsec: forwarding accountability for the SDN data plane, in: Proc. 25th Int. Conf. Comput. Commun. Netw. (ICCCN), IEEE, 2016, pp. 1--10.

	\bibitem{Bargayary2023FTISCON} % [verified; corrected from Krishnan2023]
	B. Bargayary, N. Medhi, Preserving flow table integrity in OpenFlow networks through smart contract, Cluster Computing (2023).

	\bibitem{Liu2021BFLVEC} % [venue confirmed, page range to confirm]
	H. Liu, et al., Blockchain and federated learning for collaborative intrusion detection in vehicular edge computing, IEEE Transactions on Vehicular Technology 70 (6) (2021) 6073--6084.

	\bibitem{ElDalahmeh2025} % [verified; venue corrected]
	A. El-Dalahmeh, M. Alia, Enhancing VANET security with lattice based cryptography and dynamic pseudonym updates, International Journal of Advances in Soft Computing and Its Applications 17 (2) (2025) 232--246.

	\bibitem{Wang2020Topology} % [verified]
	J. Wang, Y. Tan, J. Liu, Y. Zhang, Topology poisoning attack in SDN-enabled vehicular edge network, IEEE Internet of Things Journal 7 (10) (2020) 9563--9574.

	\bibitem{Li2022BlockREV} % [verified]
	P. Li, S. Guo, J. Wu, Q. Zhao, BlockREV: blockchain-enabled multi-controller rule enforcement verification in SDN, Security and Communication Networks 2022 (2022) 7294638.

	\bibitem{Sedjelmaci2014} % [verified]
	H. Sedjelmaci, S.~M. Senouci, M.~A. Abu-Rgheff, An efficient and lightweight intrusion detection mechanism for service-oriented vehicular networks, IEEE Internet of Things Journal 1 (6) (2014) 570--577.

	\bibitem{Sumra2011} % [verified]
	I.~A. Sumra, H. Hasbullah, I. Ahmad, Forming vehicular web of trust in VANET, in: Proc. Saudi Int. Electron. Commun. Photonics Conf. (SIECPC), IEEE, 2011.

	\bibitem{Duffield2001} % [verified]
	N.~G. Duffield, M. Grossglauser, Trajectory sampling for direct traffic observation, in: Proc. ACM SIGCOMM, 2000, pp. 271--282.

	\bibitem{Huang2020} % [verified]
	X. Huang, K. Xue, Y. Xing, D. Hu, R. Li, Q. Sun, FSDM: fast recovery saturation attack detection and mitigation framework in SDN, in: Proc. IEEE 17th Int. Conf. Mobile Ad Hoc Sensor Syst. (MASS), 2020, pp. 329--337.

	\bibitem{Chi2020} % [verified]
	P.~D. Trinh, M. Park, ECSD: enhanced compromised switch detection in an SDN-based cloud through multivariate time-series analysis, IEEE Access 8 (2020) 119346--119360.

	\bibitem{AbouElHouda2024} % [verified]
	Z. Abou El Houda, H. Moudoud, B. Brik, L. Khoukhi, Blockchain-enabled federated learning for enhanced collaborative intrusion detection in vehicular edge computing, IEEE Transactions on Intelligent Transportation Systems (2024).

	\bibitem{Mansouri2024} % [verified]
	F. Mansouri, M. Tarhouni, B. Alaya, S. Zidi, A distributed intrusion detection framework for vehicular ad hoc networks via federated learning and blockchain, Vehicular Communications (2024).

	\bibitem{Pascoal2017} % [verified]
	T.~A. Pascoal, Y.~G. Dantas, I.~E. Fonseca, V.~Nigam, Slow TCAM exhaustion DDoS attack, in: Proc. 32nd IFIP Int. Conf. ICT Syst. Secur. Privacy Prot. (IFIP SEC), Rome, Italy, 2017, pp. 17--31.

	\bibitem{Pascoal2020DoS} % [verified]
	T.~A. Pascoal, I.~E. Fonseca, V. Nigam, Slow denial-of-service attacks on software defined networks, Computer Networks 173 (2020) 107223.

	\bibitem{Cao2022LOFT} % [verified]
	J. Cao, M. Xu, Q. Li, K. Sun, Y. Yang, The LOFT attack: overflowing SDN flow tables at a low rate, IEEE/ACM Transactions on Networking 31 (2022) 1913--1928.

	\bibitem{Tang2023SFTOGuard} % [venue confirmed as this journal; full detail to confirm]
	D. Tang, W. Zhang, et al., SFTO Guard: real time detection and mitigation system for slow rate flow table overflow attacks, Journal of Network and Computer Applications (2023).

	\bibitem{Arsalan2018} % [verified]
	A. Arsalan, R.~A. Rehman, Prevention of timing attack in software defined named data network with VANETs, in: Proc. Int. Conf. Front. Inf. Technol. (FIT), Islamabad, Pakistan, 2018, pp. 247--252.

	\bibitem{Nayak2022} % [verified]
	R.~P. Nayak, S. Sethi, S.~K. Bhoi, D. Mohapatra, R.~R. Sahoo, P.~K. Sharma, D. Puthal, TFMD-SDVN: a trust framework for misbehavior detection in the edge of software defined vehicular network, Journal of Supercomputing 78 (6) (2022) 7948--7981.

	\bibitem{sae_j2735} % [verified; standard]
	SAE International, J2735: V2X Communications Message Set Dictionary, SAE International Standard, 2020.

	\bibitem{sumo_randomtrips} % [verified; tool documentation]
	German Aerospace Center (DLR), randomTrips.py, SUMO -- Simulation of Urban Mobility, https://sumo.dlr.de.

	\bibitem{FCC03-324} % [verified; regulatory order]
	Federal Communications Commission, Report and Order, ET Docket No. 03-324, Dedicated Short Range Communications in the 5.9 GHz band, 2003.

	\bibitem{pukelsheim1994three} % [verified]
	F. Pukelsheim, The three sigma rule, The American Statistician 48 (2) (1994) 88--91.

	\bibitem{Phantom2026}
	% TODO: cite as companion paper under review, add manuscript identifier once assigned.
	Anonymous, PHANTOM: post quantum mobility adaptive zero trust detection of selective time delay attacks via STARK timing attestation, ML-DSA-87 aggregate signatures, and controller failover blockchain quarantine in software defined vehicular networks, under review.

	\bibitem{Nexus2026}
	% TODO: cite as companion paper under review, add manuscript identifier once assigned.
	Anonymous, NEXUS: unified post quantum zero trust detection of selective time delay and hidden forwarding attacks via federated LSTM and controller failover blockchain in SDVN, under review.

\end{thebibliography}

\end{document}